\nonstopmode{}
\documentclass[a4paper]{book}
\usepackage[times,inconsolata,hyper]{Rd}
\usepackage{makeidx}
\makeatletter\@ifl@t@r\fmtversion{2018/04/01}{}{\usepackage[utf8]{inputenc}}\makeatother
% \usepackage{graphicx} % @USE GRAPHICX@
\makeindex{}
\begin{document}
\chapter*{}
\begin{center}
{\textbf{\huge Package `modernGLMM'}}
\par\bigskip{\large \today}
\end{center}
\ifthenelse{\boolean{Rd@use@hyper}}{\hypersetup{pdftitle = {modernGLMM: Unified R Workflows for Generalized Linear Mixed Models}}}{}
\ifthenelse{\boolean{Rd@use@hyper}}{\hypersetup{pdfauthor = {Muhammad Yaseen; Adeela Munawar; Walter W. Stroup; Kent M. Eskridge; Marina Ptukhina; Soumit Garai}}}{}
\begin{description}
\raggedright{}
\item[Type]\AsIs{Package}
\item[Title]\AsIs{Unified R Workflows for Generalized Linear Mixed Models}
\item[Version]\AsIs{2.0.0}
\item[Maintainer]\AsIs{Muhammad Yaseen }\email{myaseen208@gmail.com}\AsIs{}
\item[Description]\AsIs{Provides unified, reproducible R implementations of examples
from Stroup, W. W., Ptukhina, M., and Garai, S. (2024) ``Generalized
Linear Mixed Models: Modern Concepts, Methods and Applications'', second
edition, CRC Press, bridging the fragmented mixed-model ecosystem in R
through consistent workflows, numerical verification, and modern
statistical tooling. Companion workflows span lme4, glmmTMB, nlme,
mgcv, gamm4, brms, survival, coxme, emmeans, DHARMa, and report,
covering Gaussian and non-Gaussian responses, multi-level designs,
BLUPs, treatment structures, count data, time-to-event models,
repeated measures, spatial correlation, Bayesian inference, and power
analysis. Intended for reproducible teaching, applied GLMM analysis,
and systematic verification of model-fitting workflows.}
\item[URL]\AsIs{}\url{https://github.com/myaseen208/modernGLMM}\AsIs{}
\item[BugReports]\AsIs{}\url{https://github.com/myaseen208/modernGLMM/issues}\AsIs{}
\item[License]\AsIs{GPL-3}
\item[Encoding]\AsIs{UTF-8}
\item[Language]\AsIs{en-US}
\item[LazyData]\AsIs{true}
\item[Depends]\AsIs{R (>= 4.1.0)}
\item[Imports]\AsIs{broom.mixed, car, collapse, emmeans, ggplot2, lmerTest, MASS,
nlme, parameters, splines, stats}
\item[Suggests]\AsIs{brms, DHARMa, coxme, gamm4, glmmTMB, lme4, gratia, ggsurvfit,
insight, knitr, marginaleffects, mgcv, multcomp, nnet,
performance, plotly, quarto, report, rmarkdown, survival,
testthat (>= 3.0.0)}
\item[Config/testthat/edition]\AsIs{3}
\item[RoxygenNote]\AsIs{7.3.3}
\item[VignetteBuilder]\AsIs{quarto}
\item[NeedsCompilation]\AsIs{no}
\item[Author]\AsIs{Muhammad Yaseen [aut, cre, cph] (ORCID:
  <}\url{https://orcid.org/0000-0002-5923-1714}\AsIs{>),
Adeela Munawar [aut, ctb],
Walter W. Stroup [aut, ctb],
Kent M. Eskridge [aut, ctb],
Marina Ptukhina [aut, ctb],
Soumit Garai [aut, ctb]}
\end{description}
\Rdcontents{Contents}
\HeaderA{modernGLMM-package}{modernGLMM: Unified R Workflows for Generalized Linear Mixed Models}{modernGLMM.Rdash.package}
\aliasA{modernGLMM}{modernGLMM-package}{modernGLMM}
\keyword{internal}{modernGLMM-package}
%
\begin{Description}
`modernGLMM` provides unified, reproducible R implementations of examples
from Stroup, W. W., Ptukhina, M., and Garai, S. (2024),
\emph{Generalized Linear Mixed Models: Modern Concepts, Methods and
Applications}, second edition, CRC Press.

The 2024 book uses SAS for the worked examples and notes that the R GLMM
ecosystem is powerful but fragmented. This package responds to that gap by
collecting coherent R workflows for the book examples and by documenting how
modern R tools fit together for estimation, inference, diagnostics,
visualization, and reporting.
\end{Description}
%
\begin{Section}{Scope}

The package includes datasets, example code, vignettes, and verification
artifacts for linear models, generalized linear models, linear mixed models,
and generalized linear mixed models. Coverage spans designed experiments,
best linear unbiased prediction, treatment structures, multi-level designs,
count data, rates and proportions, multinomial responses, time-to-event data,
smoothing splines, repeated measures, correlated errors, Bayesian GLMMs, and
power analysis.
\end{Section}
%
\begin{Section}{R Ecosystem}

Workflows use explicit namespaces and draw from the R mixed-model ecosystem:
\begin{description}

\item[\pkg{lme4}] Linear and generalized linear mixed models.
\item[\pkg{glmmTMB}] Flexible GLMMs, including count and zero-inflated models.
\item[\pkg{nlme}] Correlated-error and repeated-measures models.
\item[\pkg{mgcv} and \pkg{gamm4}] Smoothing splines and additive mixed models.
\item[\pkg{brms}] Bayesian GLMM implementations.
\item[\pkg{survival} and \pkg{coxme}] Time-to-event and frailty models.
\item[\pkg{emmeans}] Estimated marginal means and contrasts.
\item[\pkg{DHARMa}] Simulation-based residual diagnostics.
\item[\pkg{report}] Readable model summaries.

\end{description}

\end{Section}
%
\begin{Section}{Workflow Example}

\begin{alltt}

data(DataSet11.1)
fit <- stats::glm(y ~ trt, family = stats::poisson(link = "log"), data = DataSet11.1)
summary(fit)

if (requireNamespace("emmeans", quietly = TRUE)) \{
  emmeans::emmeans(fit, specs = ~ trt, type = "response")
\}

if (requireNamespace("report", quietly = TRUE)) \{
  report::report(fit)
\}
\end{alltt}

\end{Section}
%
\begin{Section}{GLMM Model Framework}

A generalized linear mixed model can be written as
\deqn{g(\mu_{ij}) = \mathbf{x}_{ij}^\top \boldsymbol{\beta} + \mathbf{z}_{ij}^\top \mathbf{b}_i}{}
with random effects \eqn{\mathbf{b}_i \sim \mathcal{N}(\mathbf{0}, \mathbf{G})}{}
and conditional responses drawn from an exponential-family distribution.
\end{Section}
%
\begin{Author}
\begin{enumerate}

\item{} Muhammad Yaseen (\email{myaseen208@gmail.com})
\item{} Adeela Munawar (\email{adeela.uaf@gmail.com})
\item{} Walter W. Stroup (\email{wstroup@unl.edu})
\item{} Kent M. Eskridge (\email{keskridge1@unl.edu})

\end{enumerate}

\end{Author}
%
\begin{References}
\begin{enumerate}

\item{} Stroup, W. W., Ptukhina, M., and Garai, S. (2024).
\emph{Generalized Linear Mixed Models: Modern Concepts, Methods and
Applications} (2nd ed.). CRC Press.
\item{} Bates, D., Maechler, M., Bolker, B., and Walker, S. (2015).
Fitting linear mixed-effects models using lme4.
\emph{Journal of Statistical Software}, 67(1), 1-48.
\item{} Lenth, R. V. (2023). emmeans: Estimated Marginal Means,
aka Least-Squares Means. R package.

\end{enumerate}

\end{References}
%
\begin{SeeAlso}
\begin{itemize}

\item{} \url{https://github.com/myaseen208/modernGLMM}

\end{itemize}

\end{SeeAlso}
\HeaderA{DataExam2.B.2}{Data for Example 2.B.2 from Generalized Linear Mixed Models: Modern Concepts, Methods and Applications by Stroup, Ptukhina, and Garai (2024, 2nd ed.)}{DataExam2.B.2}
\keyword{datasets}{DataExam2.B.2}
%
\begin{Description}
Exam2.B.2 is used to visualize the effect of glm model statement with binomial data with logit and probit links.
\end{Description}
%
\begin{Usage}
\begin{verbatim}
data(DataExam2.B.2)
\end{verbatim}
\end{Usage}
%
\begin{Format}
A \code{data.frame} with 11 rows and 3 variables.
\end{Format}
%
\begin{Details}
\begin{itemize}

\item{} x independent variable
\item{} n Bernoulli trials (Bernoulli outcomes on each individual)
\item{} y number of successes on each individual

\end{itemize}

\end{Details}
%
\begin{Author}
\begin{enumerate}

\item{}  Muhammad Yaseen (\email{myaseen208@gmail.com})
\item{} Adeela Munawar (\email{adeela.uaf@gmail.com})

\end{enumerate}

\end{Author}
%
\begin{References}
\begin{enumerate}

\item{} Stroup, W. W., Ptukhina, M., and Garai, S. (2024).\emph{Generalized linear mixed models: modern concepts, methods and applications}.
CRC Press.

\end{enumerate}

\end{References}
%
\begin{SeeAlso}
\LinkA{Exam2.B.2}{Exam2.B.2}
\end{SeeAlso}
%
\begin{Examples}
\begin{ExampleCode}
data(DataExam2.B.2)
\end{ExampleCode}
\end{Examples}
\HeaderA{DataExam2.B.3}{Data for Example 2.B.3 from Generalized Linear Mixed Models: Modern Concepts, Methods and Applications by Stroup, Ptukhina, and Garai (2024, 2nd ed.)}{DataExam2.B.3}
\keyword{datasets}{DataExam2.B.3}
%
\begin{Description}
Exam2.B.3 is used to illustrate one way treatment design with Gaussian observations.
\end{Description}
%
\begin{Usage}
\begin{verbatim}
data(DataExam2.B.3)
\end{verbatim}
\end{Usage}
%
\begin{Format}
A \code{data.frame} with 6 rows and 2 variables.
\end{Format}
%
\begin{Details}
\begin{itemize}

\item{} trt treatments as factor with number 1 to 3
\item{} y   response variable

\end{itemize}

\end{Details}
%
\begin{Author}
\begin{enumerate}

\item{}  Muhammad Yaseen (\email{myaseen208@gmail.com})
\item{} Adeela Munawar (\email{adeela.uaf@gmail.com})

\end{enumerate}

\end{Author}
%
\begin{References}
\begin{enumerate}

\item{} Stroup, W. W., Ptukhina, M., and Garai, S. (2024).\emph{Generalized linear mixed models: modern concepts, methods and applications}.
CRC Press.

\end{enumerate}

\end{References}
%
\begin{SeeAlso}
\code{\LinkA{Exam2.B.3}{Exam2.B.3}}
\end{SeeAlso}
%
\begin{Examples}
\begin{ExampleCode}
data(DataExam2.B.3)
\end{ExampleCode}
\end{Examples}
\HeaderA{DataExam2.B.4}{Data for Example 2.B.4 from Generalized Linear Mixed Models: Modern Concepts, Methods and Applications by Stroup, Ptukhina, and Garai (2024, 2nd ed.)}{DataExam2.B.4}
\keyword{datasets}{DataExam2.B.4}
%
\begin{Description}
Exam2.B.4 is used to illustrate one way treatment design with Binomial observations.
\end{Description}
%
\begin{Usage}
\begin{verbatim}
data(DataExam2.B.4)
\end{verbatim}
\end{Usage}
%
\begin{Format}
A \code{data.frame} with 6 rows and 4 variables.
\end{Format}
%
\begin{Details}
\begin{itemize}

\item{} obs number of observations
\item{} trt three treatments with class factor
\item{} Nij  number of Bernoulli trials on each individual
\item{} y  number of successes on each individual

\end{itemize}

\end{Details}
%
\begin{Author}
\begin{enumerate}

\item{}  Muhammad Yaseen (\email{myaseen208@gmail.com})
\item{} Adeela Munawar (\email{adeela.uaf@gmail.com})

\end{enumerate}

\end{Author}
%
\begin{References}
\begin{enumerate}

\item{} Stroup, W. W., Ptukhina, M., and Garai, S. (2024).\emph{Generalized linear mixed models: modern concepts, methods and applications}.
CRC Press.

\end{enumerate}

\end{References}
%
\begin{SeeAlso}
\code{\LinkA{Exam2.B.4}{Exam2.B.4}}
\end{SeeAlso}
%
\begin{Examples}
\begin{ExampleCode}
data(DataExam2.B.4)
\end{ExampleCode}
\end{Examples}
\HeaderA{DataExam2.B.7}{Data for Example 2.B.7 from Generalized Linear Mixed Models: Modern Concepts, Methods and Applications by Stroup, Ptukhina, and Garai (2024, 2nd ed.)}{DataExam2.B.7}
\keyword{datasets}{DataExam2.B.7}
%
\begin{Description}
Exam2.B.7 is related to multi batch regression data assuming different forms of linear models with factorial experiment.
\end{Description}
%
\begin{Usage}
\begin{verbatim}
data(DataExam2.B.7)
\end{verbatim}
\end{Usage}
%
\begin{Format}
A \code{data.frame} with 16 rows and 4 variables.
\end{Format}
%
\begin{Details}
\begin{itemize}

\item{} Rep number of replications
\item{} a factor with two levels 1 and 2
\item{} b factor with two levels 1 and 2
\item{} y  response variable

\end{itemize}

\end{Details}
%
\begin{Author}
\begin{enumerate}

\item{}  Muhammad Yaseen (\email{myaseen208@gmail.com})
\item{} Adeela Munawar (\email{adeela.uaf@gmail.com})

\end{enumerate}

\end{Author}
%
\begin{References}
\begin{enumerate}

\item{} Stroup, W. W., Ptukhina, M., and Garai, S. (2024).\emph{Generalized linear mixed models: modern concepts, methods and applications}.
CRC Press.

\end{enumerate}

\end{References}
%
\begin{SeeAlso}
\code{\LinkA{Exam2.B.7}{Exam2.B.7}}
\end{SeeAlso}
%
\begin{Examples}
\begin{ExampleCode}
data(DataExam2.B.7)
\end{ExampleCode}
\end{Examples}
\HeaderA{DataSet10.1}{Data for Example 10.1 from Generalized Linear Mixed Models: Modern Concepts, Methods and Applications by Stroup, Ptukhina, and Garai (2024, 2nd ed.)}{DataSet10.1}
\keyword{datasets}{DataSet10.1}
%
\begin{Description}
DataSet10.1 is used for One-way random effects model and BLUP estimation
\end{Description}
%
\begin{Usage}
\begin{verbatim}
data(DataSet10.1)
\end{verbatim}
\end{Usage}
%
\begin{Format}
A \code{data.frame} with 24 rows and 2 variables.
\end{Format}
%
\begin{Details}
\begin{itemize}

\item{} a  is a factor with 12 levels
\item{} y  is a Gaussian response variable

\end{itemize}

\end{Details}
%
\begin{Author}
Muhammad Yaseen (\email{myaseen208@gmail.com})
Adeela Munawar (\email{adeela.uaf@gmail.com})
\end{Author}
%
\begin{References}
\begin{enumerate}

\item{} Stroup, W. W., Ptukhina, M., and Garai, S. (2024).
\emph{Generalized Linear Mixed Models: Modern Concepts, Methods and Applications (2nd ed.)}.
CRC Press.

\end{enumerate}

\end{References}
%
\begin{SeeAlso}
\code{\LinkA{Exam10.1}{Exam10.1}}
\end{SeeAlso}
%
\begin{Examples}
\begin{ExampleCode}
data(DataSet10.1)
\end{ExampleCode}
\end{Examples}
\HeaderA{DataSet10.2}{Data for Example 10.2 from Generalized Linear Mixed Models: Modern Concepts, Methods and Applications by Stroup, Ptukhina, and Garai (2024, 2nd ed.)}{DataSet10.2}
\keyword{datasets}{DataSet10.2}
%
\begin{Description}
DataSet10.2 is used for Two-way nested random effects model and BLUP estimation
\end{Description}
%
\begin{Usage}
\begin{verbatim}
data(DataSet10.2)
\end{verbatim}
\end{Usage}
%
\begin{Format}
A \code{data.frame} with 28 rows and 3 variables with levels of b nested within levels of a.
\end{Format}
%
\begin{Details}
\begin{itemize}

\item{} a  is a factor with 7 levels
\item{} b  is a factor with 2 levels
\item{} y  is a Gaussian response variable

\end{itemize}

\end{Details}
%
\begin{Author}
Muhammad Yaseen (\email{myaseen208@gmail.com})
Adeela Munawar (\email{adeela.uaf@gmail.com})
\end{Author}
%
\begin{References}
\begin{enumerate}

\item{} Stroup, W. W., Ptukhina, M., and Garai, S. (2024).
\emph{Generalized Linear Mixed Models: Modern Concepts, Methods and Applications (2nd ed.)}.
CRC Press.

\end{enumerate}

\end{References}
%
\begin{SeeAlso}
\code{\LinkA{Exam10.2}{Exam10.2}}
\end{SeeAlso}
%
\begin{Examples}
\begin{ExampleCode}
data(DataSet10.2)
\end{ExampleCode}
\end{Examples}
\HeaderA{DataSet10.4}{Data for Example 10.4 from Generalized Linear Mixed Models: Modern Concepts, Methods and Applications by Stroup, Ptukhina, and Garai (2024, 2nd ed.)}{DataSet10.4}
\keyword{datasets}{DataSet10.4}
%
\begin{Description}
DataSet10.4 is used for demonstrating the relationship between BLUP and Fixed Effect Estimators
\end{Description}
%
\begin{Usage}
\begin{verbatim}
data(DataSet10.4)
\end{verbatim}
\end{Usage}
%
\begin{Format}
A \code{data.frame} with 32 rows and 3 variables.
\end{Format}
%
\begin{Details}
\begin{itemize}

\item{} a  is a factor with 2 levels
\item{} b  is a factor with 8 levels
\item{} y  is a Gaussian response variable

\end{itemize}

\end{Details}
%
\begin{Author}
Muhammad Yaseen (\email{myaseen208@gmail.com})
Adeela Munawar (\email{adeela.uaf@gmail.com})
\end{Author}
%
\begin{References}
\begin{enumerate}

\item{} Stroup, W. W., Ptukhina, M., and Garai, S. (2024).
\emph{Generalized Linear Mixed Models: Modern Concepts, Methods and Applications (2nd ed.)}.
CRC Press.

\end{enumerate}

\end{References}
%
\begin{SeeAlso}
\code{\LinkA{Exam10.4}{Exam10.4}}
\end{SeeAlso}
%
\begin{Examples}
\begin{ExampleCode}
data(DataSet10.4)
\end{ExampleCode}
\end{Examples}
\HeaderA{DataSet11.1}{Data Set 11.1 - Completely Randomized Count Data}{DataSet11.1}
%
\begin{Description}
Two-treatment completely randomized count data used in Chapter 11,
Section 11.1.2 of Stroup, Ptukhina, and Garai (2024) to compare
pre-GLM ANOVA analyses, a Poisson GLM, and a Poisson-normal GLMM with
an observation-level random effect.
\end{Description}
%
\begin{Usage}
\begin{verbatim}
data(DataSet11.1)
\end{verbatim}
\end{Usage}
%
\begin{Format}
A \code{data.frame} with 10 rows and 3 variables:
\begin{description}

\item[trt] Treatment factor with levels \code{1} and \code{2}.
\item[unit] Replication unit within treatment.
\item[count] Non-negative integer count response.

\end{description}

\end{Format}
%
\begin{Details}
The reconstructed counts are \code{2, 2, 3, 5, 12} for treatment 1 and
\code{6, 11, 12, 12, 34} for treatment 2.

**Data provenance:** Reconstructed.

**Reconstruction method:** The integer counts were solved from the
published treatment means, log-count means, ANOVA residual mean square,
and log-count ANOVA residual mean square in Chapter 11, Section 11.1.2.

**Book agreement:** Exact for the published untransformed ANOVA treatment
means, F statistic, p-value, log-count treatment means, log-count standard
error, and log-count treatment comparison. Approximate for the Poisson-normal
GLMM when fitted in R because \pkg{lme4} uses likelihood-based estimation,
whereas the book reports SAS PROC GLIMMIX output.

**Limitations:** The book references the SAS file \code{Data\_Set\_11\_1},
but the raw SAS file is not included in this package or in the uploaded PDF.
The reconstruction is therefore based on printed numerical constraints.
\end{Details}
%
\begin{References}
Stroup, W. W., Ptukhina, M., and Garai, S. (2024).
\emph{Generalized Linear Mixed Models: Modern Concepts, Methods and
Applications}, 2nd ed. CRC Press.
\end{References}
%
\begin{SeeAlso}
\code{\LinkA{Exam11.1}{Exam11.1}}
\end{SeeAlso}
%
\begin{Examples}
\begin{ExampleCode}
data(DataSet11.1)
stats::aggregate(count ~ trt, data = DataSet11.1, FUN = mean)
\end{ExampleCode}
\end{Examples}
\HeaderA{DataSet11.3}{Data Set 11.3 - Randomized Complete Block Count Data}{DataSet11.3}
%
\begin{Description}
Three-treatment randomized complete block count data used in Chapter 11,
Section 11.4 of Stroup, Ptukhina, and Garai (2024) to illustrate
overdispersion in blocked count-data models and alternatives including
marginal Poisson, Poisson-normal, negative binomial, and generalized
Poisson models.
\end{Description}
%
\begin{Usage}
\begin{verbatim}
data(DataSet11.3)
\end{verbatim}
\end{Usage}
%
\begin{Format}
A \code{data.frame} with 30 rows and 3 variables:
\begin{description}

\item[block] Block factor with 10 levels.
\item[trt] Treatment factor with 3 levels.
\item[count] Non-negative integer count response.

\end{description}

\end{Format}
%
\begin{Details}
Treatment sample means are exactly \code{8.0}, \code{13.3}, and
\code{25.7}, matching the printed marginal means in Section 11.4.1.

**Data provenance:** Reconstructed.

**Reconstruction method:** Counts were reconstructed from the published
treatment sample means and selected model summaries in Chapter 11,
Section 11.4. The optimization target used the printed naive Poisson,
Poisson-normal, and negative-binomial fitted means, Pearson chi-square/DF,
and variance-component summaries as calibration constraints.

**Book agreement:** Exact for treatment sample means. Approximate for
R-fitted Poisson and negative-binomial GLMM summaries. The printed SAS
working-covariance marginal model and generalized Poisson fit are not
directly reproducible with base package dependencies.

**Limitations:** The original SAS data file was not available in the
workspace, the uploaded PDF, or public publisher materials located during
verification. Several integer datasets satisfy the printed summaries; this
file is a constrained reconstruction, not a confirmed original data file.
\end{Details}
%
\begin{References}
Stroup, W. W., Ptukhina, M., and Garai, S. (2024).
\emph{Generalized Linear Mixed Models: Modern Concepts, Methods and
Applications}, 2nd ed. CRC Press.
\end{References}
%
\begin{SeeAlso}
\code{\LinkA{Exam11.3}{Exam11.3}}
\end{SeeAlso}
%
\begin{Examples}
\begin{ExampleCode}
data(DataSet11.3)
stats::aggregate(count ~ trt, data = DataSet11.3, FUN = mean)
\end{ExampleCode}
\end{Examples}
\HeaderA{DataSet11.4}{Data Set 11.4 - Split-Plot Count Data}{DataSet11.4}
%
\begin{Description}
Split-plot count data for the multi-level Chapter 11 example in
Section 11.5 of Stroup, Ptukhina, and Garai (2024). The treatment
structure is a 7 by 4 factorial with four blocks.
\end{Description}
%
\begin{Usage}
\begin{verbatim}
data(DataSet11.4)
\end{verbatim}
\end{Usage}
%
\begin{Format}
A \code{data.frame} with 112 rows and 4 variables:
\begin{description}

\item[block] Block factor with 4 levels.
\item[a] Whole-plot treatment factor with 7 levels.
\item[b] Subplot treatment factor with 4 levels.
\item[count] Non-negative integer count response.

\end{description}

\end{Format}
%
\begin{Details}
The data provide a complete 4-block, 7-level \code{a}, 4-level \code{b}
split-plot layout for fitting the marginal Poisson and negative-binomial
GLMMs described in Chapter 11, Section 11.5.

**Data provenance:** Synthetic.

**Reconstruction method:** The original raw SAS data were not available.
Counts were generated from the printed model-scale and data-scale means
for the displayed \code{a = 1} and \code{a = 7} cells, with intermediate
\code{a} levels interpolated to preserve the 7 by 4 factorial structure.

**Book agreement:** Approximate for the displayed \code{a = 1} and
\code{a = 7} fitted mean pattern. Not verified for all 28 cells because
the book prints only selected least-squares means.

**Limitations:** This dataset is suitable for runnable examples and
demonstrations of the Chapter 11 workflow, but it is not a verified copy
of the book's original \code{sp\_counts} data.
\end{Details}
%
\begin{References}
Stroup, W. W., Ptukhina, M., and Garai, S. (2024).
\emph{Generalized Linear Mixed Models: Modern Concepts, Methods and
Applications}, 2nd ed. CRC Press.
\end{References}
%
\begin{SeeAlso}
\code{\LinkA{Exam11.4}{Exam11.4}}
\end{SeeAlso}
%
\begin{Examples}
\begin{ExampleCode}
data(DataSet11.4)
stats::aggregate(count ~ a + b, data = DataSet11.4, FUN = mean)
\end{ExampleCode}
\end{Examples}
\HeaderA{DataSet12.1}{Data for Example 12.1 — Continuous Proportion Dose-Response (Chapter 12)}{DataSet12.1}
\keyword{datasets}{DataSet12.1}
%
\begin{Description}
Continuous proportion data from a dose-response study with two treatments
each applied to 12 runs, observed at six dose levels (0–5).  Implements
Section 12.6.2 of Stroup et al. (2024), Data Set 12.7 (SAS name
\code{rates}).  The response is a continuous proportion suitable for
beta regression with a run-level GLMM random effect.

\strong{Reconstruction note:} The actual data appear in the SAS Data and
Program Library as Data Set 12.7.  This version is reconstructed from the
published regression parameters (\eqn{\beta_{0,0}=0.6965}{},
\eqn{\beta_{1,0}=0.2846}{}; \eqn{\beta_{0,1}=0.8054}{},
\eqn{\beta_{1,1}=0.5541}{}; p.406) with run random effects drawn from
\eqn{N(0, 0.45^2)}{} and Beta precision \eqn{\phi=10}{}, using
\code{set.seed(2024)}.  Numerical output will approximate published values.
\end{Description}
%
\begin{Usage}
\begin{verbatim}
data(DataSet12.1)
\end{verbatim}
\end{Usage}
%
\begin{Format}
A \code{data.frame} with 144 rows and 4 variables:
\begin{description}

\item[trt] Treatment factor with levels 0 and 1
\item[run] Run factor with 24 levels (t0r1-t0r12, t1r1-t1r12);
12 runs per treatment, nested within treatment
\item[dose] Integer dose level: 0, 1, 2, 3, 4, or 5
\item[proportion] Continuous proportion response in \eqn{(0, 1)}{}

\end{description}

\end{Format}
%
\begin{Details}
The beta GLMM for the linear dose-response within each treatment is:
\deqn{\text{logit}(\mu_{ijk}) = \beta_{0i} + \beta_{1i} D_j + r(\tau)_{ik}}{}
where \eqn{r(\tau)_{ik} \sim \mathcal{N}(0,\sigma^2_r)}{} is the run random
effect nested within treatment, and \eqn{D_j}{} is the dose level.
\end{Details}
%
\begin{References}
\begin{enumerate}

\item{} Stroup, W. W., Ptukhina, M., and Garai, S. (2024).
\emph{Generalized Linear Mixed Models: Modern Concepts,
Methods and Applications}. CRC Press. Section 12.6.2, pp.404-407.
\item{} Ferrari, S., \& Cribari-Neto, F. (2004).
Beta regression for modelling rates and proportions.
\emph{Journal of Applied Statistics}, 31(7), 799-815.

\end{enumerate}

\end{References}
%
\begin{SeeAlso}
\code{\LinkA{Exam12.1}{Exam12.1}}
\end{SeeAlso}
%
\begin{Examples}
\begin{ExampleCode}
data(DataSet12.1)
str(DataSet12.1)
with(DataSet12.1, tapply(proportion, list(trt, dose), mean))
\end{ExampleCode}
\end{Examples}
\HeaderA{DataSet12.2}{Data for Example 12.2 — Binomial Nested Factorial (Chapter 12)}{DataSet12.2}
\keyword{datasets}{DataSet12.2}
%
\begin{Description}
Binomial count data from a nested factorial design with factor A (2 levels,
representing two sets) and factor B (3 levels nested within A), replicated
across 10 blocks (5 blocks per set).  Implements Section 12.3.2 of
Stroup et al. (2024), Data Set 12.2 (SAS name \code{ch11\_ex\_12C2}).
The response is the number of successes out of a fixed total, suitable for
binomial GLMM analysis.

\strong{Reconstruction note:} The actual data appear in the SAS Data and
Program Library as Data Set 12.2.  This version is reconstructed from the
published design description (5 blocks × setA + 5 blocks × setB, 3
treatments B0/B1/B2 per block, N=20 per cell, 30 obs total) and the
published logit B(A) LSMeans (p.382) using \code{set.seed(123)}.
Numerical output will approximate published values.
\end{Description}
%
\begin{Usage}
\begin{verbatim}
data(DataSet12.2)
\end{verbatim}
\end{Usage}
%
\begin{Format}
A \code{data.frame} with 30 rows and 5 variables:
\begin{description}

\item[block] Block factor with 10 levels (1–10); blocks 1–5 are
assigned to \code{setA}, blocks 6–10 to \code{setB}
\item[a] Set factor with 2 levels (\code{setA}, \code{setB});
each set assigned to 5 blocks (whole-plot factor)
\item[b] Treatment factor nested within A, with 3 levels
(\code{B0}, \code{B1}, \code{B2}) within each set
\item[f] Number of successes (integer count)
\item[n] Total number of trials per cell (integer, N=20)

\end{description}

\end{Format}
%
\begin{Details}
The binomial GLMM for the nested factorial structure is:
\deqn{\text{logit}(\pi_{ijk}) = \eta + \alpha_i + \beta_{j(i)} + r_{k(i)}}{}
where \eqn{\alpha_i}{} is the set (A) effect, \eqn{\beta_{j(i)}}{} is the
treatment B nested within A, and \eqn{r_{k(i)} \sim
\mathcal{N}(0, \sigma^2_A)}{} is the block-within-A random effect
(whole-plot error).

Published results (2nd ed. p.382):
\eqn{\hat\sigma^2_A = 0.4784}{} (SE=0.2664); A: F(1,8)=0.46, p=0.5187;
B(A): F(4,16)=2.43, p=0.0907.
\end{Details}
%
\begin{References}
\begin{enumerate}

\item{} Stroup, W. W., Ptukhina, M., and Garai, S. (2024).
\emph{Generalized Linear Mixed Models: Modern Concepts,
Methods and Applications}. CRC Press. Section 12.3.2, pp.388-392.

\end{enumerate}

\end{References}
%
\begin{SeeAlso}
\code{\LinkA{Exam12.2}{Exam12.2}}
\end{SeeAlso}
%
\begin{Examples}
\begin{ExampleCode}
data(DataSet12.2)
str(DataSet12.2)
with(DataSet12.2, tapply(f / n, list(a, b), mean))
\end{ExampleCode}
\end{Examples}
\HeaderA{DataSet14.1}{Data for Example 14.1 — Ordinal Proportional Odds GLMM (Chapter 14)}{DataSet14.1}
\keyword{datasets}{DataSet14.1}
%
\begin{Description}
Ordinal frequency count data from a randomised complete block design with
6 treatments and 10 blocks.  The response is a three-category ordered rating
(slight, moderate, severe).  Implements Section 14.3 of Stroup et al. (2024),
Data Set 14.1 (SAS name \code{univar\_multinom}).  Suitable for proportional-odds
(cumulative logit) GLMM analysis.

\strong{Reconstruction note:} The first 9 rows (block 1, treatments 0–2)
are transcribed exactly from the published table (p.435 of the 2nd ed.).
The remaining 171 rows are reconstructed from the published cumulative-logit
parameters (\eqn{\hat\eta_0=0.3492}{}, \eqn{\hat\eta_1=1.9956}{}, treatment
effects pp.436-437) using \code{set.seed(2024)}.  The actual data appear in
the SAS Data and Program Library as Data Set 14.1.
\end{Description}
%
\begin{Usage}
\begin{verbatim}
data(DataSet14.1)
\end{verbatim}
\end{Usage}
%
\begin{Format}
A \code{data.frame} with 180 rows and 4 variables:
\begin{description}

\item[blk] Block factor with 10 levels (1–10)
\item[trt] Treatment factor with 6 levels (0–5)
\item[rating] Ordered factor with levels \code{slight} < \code{modrat} <
\code{severe}
\item[y] Frequency count of observations in this blk × trt × rating cell
(non-negative integer)

\end{description}

\end{Format}
%
\begin{Details}
The proportional-odds GLMM for treatment \eqn{i}{} in block \eqn{k}{} with
ordinal category \eqn{j}{} is:
\deqn{\text{logit}[P(Y_{ik} \le j)] = \eta_j - \tau_i - b_k}{}
where \eqn{\eta_j}{} are the cumulative intercept thresholds,
\eqn{\tau_i}{} is the treatment fixed effect, and
\eqn{b_k \sim \mathcal{N}(0, \sigma^2_b)}{} is the block random effect.

Published results (2nd ed. pp.436-437):
\eqn{\hat\eta_0=0.3492}{}, \eqn{\hat\eta_1=1.9956}{};
Treatment F(5,768)=17.67, p<0.0001.
\end{Details}
%
\begin{References}
\begin{enumerate}

\item{} Stroup, W. W., Ptukhina, M., and Garai, S. (2024).
\emph{Generalized Linear Mixed Models: Modern Concepts,
Methods and Applications}. CRC Press. Section 14.3, pp.434-438.
\item{} Agresti, A. (2010). \emph{Analysis of Ordinal Categorical Data},
2nd ed. Wiley.

\end{enumerate}

\end{References}
%
\begin{SeeAlso}
\code{\LinkA{Exam14.1}{Exam14.1}}
\end{SeeAlso}
%
\begin{Examples}
\begin{ExampleCode}
data(DataSet14.1)
str(DataSet14.1)
with(DataSet14.1, tapply(y, list(trt, rating), sum))
\end{ExampleCode}
\end{Examples}
\HeaderA{DataSet14.2}{Data for Example 14.2 — Non-Proportional Odds: Poinsettia Trial (Chapter 14)}{DataSet14.2}
\keyword{datasets}{DataSet14.2}
%
\begin{Description}
Ordinal frequency count data from a poinsettia variety trial with 3 varieties
evaluated by 12 growers on a three-category ordered quality rating.  Implements
Section 14.4 of Stroup et al. (2024), Data Set 14.2 (SAS name
\code{poinsettia}).  The proportional-odds assumption fails (non-proportional
odds), requiring separate logistic regressions per category boundary.

\strong{Reconstruction note:} The marginal frequency table
(variety × rating totals across growers) is transcribed exactly from the
published table (p.438 of the 2nd ed.).  Counts are distributed
proportionally across 12 growers using \code{set.seed(2024)}.  The actual
data appear in the SAS Data and Program Library as Data Set 14.2.
\end{Description}
%
\begin{Usage}
\begin{verbatim}
data(DataSet14.2)
\end{verbatim}
\end{Usage}
%
\begin{Format}
A \code{data.frame} with 108 rows and 4 variables:
\begin{description}

\item[grower] Grower identifier (factor with 12 levels); serves as the
GLMM random-effect block
\item[variety] Poinsettia variety factor with 3 levels (1, 2, 3)
\item[rating] Quality rating factor with levels A, B, C (where A = lowest
quality category and C = highest quality category)
\item[y] Frequency count of plants in this grower × variety × rating cell
(non-negative integer)

\end{description}

\end{Format}
%
\begin{Details}
The proportional-odds (cumulative logit) model assumes a single variety
effect across all category boundaries.  The non-proportional-odds model
allows variety effects to differ by boundary \eqn{j}{}:
\deqn{\text{logit}[P(Y_{ik} \le j)] = \eta_j - \tau_{ij} - g_k}{}
where \eqn{\tau_{ij}}{} varies with \eqn{j}{} and
\eqn{g_k \sim \mathcal{N}(0, \sigma^2_g)}{} is the grower random effect.

Published marginal totals (variety × rating, p.438):

\Tabular{lrrr}{
& A & B & C \\{}
v1 & 192 & 174 & 176 \\{}
v2 & 65 & 395 & 68 \\{}
v3 & 262 & 30 & 244
}
Published results: proportional odds variety F(2,8)≈0.02, p≈0.98 (fail);
non-proportional odds F(4,8)=69.65.
\end{Details}
%
\begin{References}
\begin{enumerate}

\item{} Stroup, W. W., Ptukhina, M., and Garai, S. (2024).
\emph{Generalized Linear Mixed Models: Modern Concepts,
Methods and Applications}. CRC Press. Section 14.4, pp.438-447.

\end{enumerate}

\end{References}
%
\begin{SeeAlso}
\code{\LinkA{Exam14.2}{Exam14.2}}
\end{SeeAlso}
%
\begin{Examples}
\begin{ExampleCode}
data(DataSet14.2)
str(DataSet14.2)
## Marginal totals match published table (book p.438)
with(DataSet14.2, tapply(y, list(variety, rating), sum))
\end{ExampleCode}
\end{Examples}
\HeaderA{DataSet17.1}{Data for Example 17.1 — Crossover with ARH(1) Covariance (Chapter 17)}{DataSet17.1}
\keyword{datasets}{DataSet17.1}
%
\begin{Description}
Longitudinal data from a two-treatment crossover study with 41 subjects
measured at 6 equally-spaced time points per period (2 periods), plus a
baseline covariate.  Implements Example 17.3.1 of Stroup et al. (2024),
Data Set 17.1 (SAS name \code{x\_over\_ante1}).  Demonstrates covariance
model selection (CS, AR(1), ARH(1)) for repeated measures within a
crossover design.

\strong{Reconstruction note:} The actual data appear in the SAS Data and
Program Library as Data Set 17.1.  This version is reconstructed from the
published design description (41 subjects: 17 in sequence 0→1, 24 in
sequence 1→0; 2 periods × 6 times; baseline covariate) and published
regression contrasts (pp.513-516) using \code{set.seed(2024)}.
Published AICC and F-statistics will not be reproduced exactly.
\end{Description}
%
\begin{Usage}
\begin{verbatim}
data(DataSet17.1)
\end{verbatim}
\end{Usage}
%
\begin{Format}
A \code{data.frame} with 492 rows and 7 variables:
\begin{description}

\item[id] Subject identifier (factor with 41 levels)
\item[sequence] Crossover sequence: \code{"01"} (trt 0 then 1) or
\code{"10"} (trt 1 then 0); 17 and 24 subjects respectively
\item[period] Period factor with 2 levels (1, 2)
\item[trt] Treatment factor with 2 levels (0, 1)
\item[t] Time index within period: 0, 1, 2, 3, 4, 5 (equally spaced)
\item[baseline] Subject-level baseline covariate (continuous)
\item[y] Continuous Gaussian response

\end{description}

\end{Format}
%
\begin{Details}
The random coefficient model with ARH(1) within-subject covariance is:
\deqn{y_{ijkl} = \beta_{0i} + b_{0ik} + (\beta_{1i} + b_{1ik}) T_j +
      \rho_p + \beta_b X_k + e_{ijkl}}{}
where \eqn{b_{0ik}}{} and \eqn{b_{1ik}}{} are subject-specific random
intercept and slope, \eqn{\rho_p}{} is the period effect, \eqn{X_k}{} is
the baseline covariate, and \eqn{e_{ijkl}}{} follows an ARH(1) structure
(heterogeneous variances per time point with AR(1) lag-1 correlation).

Published results (2nd ed. p.516):
Best covariance = ARH(1), AICC=3035.25.
Equal intercepts: F(1,526)=1.16, p=0.2818;
Equal slopes: F(1,104)=4.18, p=0.0434.
\end{Details}
%
\begin{References}
\begin{enumerate}

\item{} Stroup, W. W., Ptukhina, M., and Garai, S. (2024).
\emph{Generalized Linear Mixed Models: Modern Concepts,
Methods and Applications}. CRC Press. Section 17.3.1, pp.513-516.
\item{} Verbeke, G., \& Molenberghs, G. (2000).
\emph{Linear Mixed Models for Longitudinal Data}. Springer.

\end{enumerate}

\end{References}
%
\begin{SeeAlso}
\code{\LinkA{Exam17.1}{Exam17.1}}
\end{SeeAlso}
%
\begin{Examples}
\begin{ExampleCode}
data(DataSet17.1)
str(DataSet17.1)
cat("Subjects per sequence:\n")
print(table(unique(DataSet17.1[, c("id", "sequence")])$sequence))
\end{ExampleCode}
\end{Examples}
\HeaderA{DataSet17.2}{Data for Example 17.2 — Sparse Longitudinal Data with SP(POW) (Chapter 17)}{DataSet17.2}
\keyword{datasets}{DataSet17.2}
%
\begin{Description}
Sparse longitudinal data from 41 subjects (17 on treatment 1, 24 on
treatment 2) observed at up to 9 unequally-spaced time points
(0, 1, 2, 4, 8, 16, 32, 64, 128), with only 101 of 369 possible
observations present.  Implements Example 17.3.2 of Stroup et al. (2024),
Data Set 17.2 (SAS name \code{all}).  Demonstrates the SP(POW) spatial-power
covariance as a generalisation of AR(1) to irregular time grids.

\strong{Reconstruction note:} The actual data appear in the SAS Data and
Program Library as Data Set 17.2.  This version is reconstructed from the
published design description (41 subjects, 9 unequal times, sparse 101/369
observations) and published random coefficient model parameters (pp.516-518)
using \code{set.seed(2024)}.  Published AICC and parameter estimates will
not be reproduced exactly.
\end{Description}
%
\begin{Usage}
\begin{verbatim}
data(DataSet17.2)
\end{verbatim}
\end{Usage}
%
\begin{Format}
A \code{data.frame} with 101 rows and 4 variables:
\begin{description}

\item[subject] Subject identifier (factor with 41 levels)
\item[trt] Treatment factor with 2 levels (1, 2);
17 subjects on trt=1, 24 subjects on trt=2
\item[time] Actual time of measurement: one of 0, 1, 2, 4, 8, 16, 32,
64, 128 (continuous; unequally spaced)
\item[y] Continuous Gaussian response

\end{description}

\end{Format}
%
\begin{Details}
For unequally-spaced observations, the SP(POW) covariance between
observations at times \eqn{X_j}{} and \eqn{X_{j'}}{} on the same subject is:
\deqn{\text{Cov}(e_j, e_{j'}) = \sigma^2 \rho^{|X_j - X_{j'|}}}{}
Unlike AR(1), SP(POW) accommodates unequal spacing and missing observations.
The random coefficient model is:
\deqn{y_{ik} = \beta_{0i} + b_{0k} + (\beta_{1i} + b_{1k}) X_j + e_{ijk}}{}
where \eqn{b_{0k}}{} and \eqn{b_{1k}}{} are random intercept and slope.

Published results (2nd ed. pp.517-518):
Best = heterogeneous SP(POW) by treatment, AICC=575.96;
Homogeneity test: \eqn{\chi^2}{}=8.65, p=0.0132.
\end{Details}
%
\begin{References}
\begin{enumerate}

\item{} Stroup, W. W., Ptukhina, M., and Garai, S. (2024).
\emph{Generalized Linear Mixed Models: Modern Concepts,
Methods and Applications}. CRC Press. Section 17.3.2, pp.516-518.
\item{} Littell, R.C., Milliken, G.A., Stroup, W.W., Wolfinger, R.D.,
\& Schabenberger, O. (2006). \emph{SAS for Mixed Models}, 2nd ed.
SAS Institute.

\end{enumerate}

\end{References}
%
\begin{SeeAlso}
\code{\LinkA{Exam17.2}{Exam17.2}}
\end{SeeAlso}
%
\begin{Examples}
\begin{ExampleCode}
data(DataSet17.2)
str(DataSet17.2)
cat("Observations per treatment:\n")
print(table(DataSet17.2$trt))
cat("Sparsity:", nrow(DataSet17.2), "of",
    length(unique(DataSet17.2$subject)) * 9, "possible\n")
\end{ExampleCode}
\end{Examples}
\HeaderA{DataSet18.1}{Data for Example 18.1 — Alliance Wheat Trial: Gaussian Spatial (Chapter 18)}{DataSet18.1}
\keyword{datasets}{DataSet18.1}
%
\begin{Description}
Wheat yield data from a 12×12 field grid (144 plots) with 48 treatments
and 3 complete blocks (columns 1–4, 5–8, 9–12).  Implements Section 18.3
of Stroup et al. (2024), Data Set 18.1 (SAS names \code{ch15\_ex1} /
\code{asademo}).  Demonstrates spatial covariance model selection
(spherical, exponential, Gaussian) via AICC for a Gaussian response.

\strong{Reconstruction note:} The actual data appear in the SAS Data and
Program Library as Data Set 18.1 (Alliance wheat trial).  This version is
reconstructed from the published design description (48 treatments, 12×12
grid, 3 column blocks) and published spatial parameters (range=4.1214,
\eqn{\sigma^2}{}=14.0107, pp.569-571) using \code{set.seed(2024)}.
Published AICC values and range estimates will not be reproduced exactly.
\end{Description}
%
\begin{Usage}
\begin{verbatim}
data(DataSet18.1)
\end{verbatim}
\end{Usage}
%
\begin{Format}
A \code{data.frame} with 144 rows and 5 variables:
\begin{description}

\item[row] Field row position (integer, 1–12)
\item[col] Field column position (integer, 1–12)
\item[block] Complete block factor: columns 1-4 = block 1,
columns 5-8 = block 2, columns 9-12 = block 3
\item[trt] Treatment factor with 48 levels (t01–t48);
each treatment appears once per block
\item[y] Continuous wheat yield response

\end{description}

\end{Format}
%
\begin{Details}
Spatial models replace the iid residual assumption with a covariance
function of Euclidean distance \eqn{d_{ij}}{} between plots.  The spherical
covariance (SP(SPH)), which is the best model in the book, is:
\deqn{C(d_{ij}) = \begin{cases}
  \sigma^2 \left[1 - \frac{3d_{ij}}{2r} + \frac{d_{ij}^3}{2r^3}\right]
  & d_{ij} < r \\
  0 & d_{ij} \ge r
\end{cases}}{}
where \eqn{r}{} is the range parameter and \eqn{\sigma^2}{} is the sill.

Published results (2nd ed. pp.569-571):
SP(SPH): range=4.1214, residual \eqn{\sigma^2}{}=14.0107.
AICC: Spherical=512.9, Exponential=521.9, Gaussian=530.9,
Incomplete block=588.8, Complete block=597.9.
\end{Details}
%
\begin{References}
\begin{enumerate}

\item{} Stroup, W. W., Ptukhina, M., and Garai, S. (2024).
\emph{Generalized Linear Mixed Models: Modern Concepts,
Methods and Applications}. CRC Press. Section 18.3, pp.563-573.
\item{} Pinheiro, J., \& Bates, D. (2000).
\emph{Mixed-Effects Models in S and S-PLUS}. Springer.

\end{enumerate}

\end{References}
%
\begin{SeeAlso}
\code{\LinkA{Exam18.1}{Exam18.1}}
\end{SeeAlso}
%
\begin{Examples}
\begin{ExampleCode}
data(DataSet18.1)
str(DataSet18.1)
with(DataSet18.1, table(block))
\end{ExampleCode}
\end{Examples}
\HeaderA{DataSet18.2}{Data for Example 18.2 — Hessian Fly Trial: Binomial Spatial (Chapter 18)}{DataSet18.2}
\keyword{datasets}{DataSet18.2}
%
\begin{Description}
Binomial resistance data from a 4×4 lattice trial of 16 wheat varieties
on a 12×12 field grid, with 64 plots.  Implements Section 18.4 of Stroup
et al. (2024), Data Set 18.2 (SAS name \code{HessianFly}).  Demonstrates
spatial binomial GLMM analysis with model comparison across RCB, incomplete
block, and spatial covariance structures via AICC.

\strong{Reconstruction note:} The actual data appear in the SAS Data and
Program Library as Data Set 18.2.  This version is reconstructed from the
published design description (16 varieties, 4×4 lattice, 4 incomplete
blocks, 64 plots, binomial response) and published spatial parameters
(SP(SPH)=3.2256, \eqn{\sigma^2}{}=0.5111, pp.573-579) using
\code{set.seed(2024)}.  Published AICC values will not be reproduced exactly.
\end{Description}
%
\begin{Usage}
\begin{verbatim}
data(DataSet18.2)
\end{verbatim}
\end{Usage}
%
\begin{Format}
A \code{data.frame} with 64 rows and 6 variables:
\begin{description}

\item[block] Complete block (replicate) factor with 4 levels;
each block contains all 16 varieties
\item[variety] Wheat variety factor with 16 levels (v01–v16)
\item[row] Plot row position in the 12×12 field grid (integer)
\item[col] Plot column position in the 12×12 field grid (integer)
\item[y] Number of Hessian fly-damaged plants (integer count)
\item[n] Total number of plants per plot (integer)

\end{description}

\end{Format}
%
\begin{Details}
The G-side binomial GLMM with spherical spatial random effects is:
\deqn{y_{ij} | b_{ij} \sim \text{Binomial}(n_{ij}, \pi_{ij})}{}
\deqn{\text{logit}(\pi_{ij}) = \eta + \tau_v + b_{ij}}{}
where \eqn{\tau_v}{} is the variety fixed effect and \eqn{b_{ij}}{} is a
spatial random effect with spherical covariance structure:
\deqn{\text{Cov}(b_{ij}, b_{i'j'}) = \sigma^2 \text{sph}(d_{ij,i'j'}/r)}{}

Published results (2nd ed. pp.573-579):
AICC: RCB=317.27, Incomplete block=296.64; F\_entry(RCB)=6.81.
G-side spherical: SP(SPH)=3.2256, \eqn{\sigma^2}{}=0.5111, AICC=301.91.
\end{Details}
%
\begin{References}
\begin{enumerate}

\item{} Stroup, W. W., Ptukhina, M., and Garai, S. (2024).
\emph{Generalized Linear Mixed Models: Modern Concepts,
Methods and Applications}. CRC Press. Section 18.4, pp.573-579.

\end{enumerate}

\end{References}
%
\begin{SeeAlso}
\code{\LinkA{Exam18.2}{Exam18.2}}
\end{SeeAlso}
%
\begin{Examples}
\begin{ExampleCode}
data(DataSet18.2)
str(DataSet18.2)
with(DataSet18.2, tapply(y / n, variety, mean))
\end{ExampleCode}
\end{Examples}
\HeaderA{DataSet21.1}{Data for Example 21.1 — Power Curve for One-Way ANOVA (Chapter 21)}{DataSet21.1}
\keyword{datasets}{DataSet21.1}
%
\begin{Description}
Pre-computed power values for a balanced one-way ANOVA (3 treatments)
across a grid of sample sizes and standardised effect sizes.
Suitable for plotting power curves and selecting sample sizes.
\end{Description}
%
\begin{Usage}
\begin{verbatim}
data(DataSet21.1)
\end{verbatim}
\end{Usage}
%
\begin{Format}
A \code{data.frame} with 30 rows and 3 variables:
\begin{description}

\item[n\_per\_group] Sample size per treatment group (5–50, step 5)
\item[effect\_size] Cohen's \eqn{f}{} effect size (0.2, 0.5, or 0.8)
\item[power] Estimated power (0–1)

\end{description}

\end{Format}
%
\begin{Details}
Power for a one-way ANOVA F-test is computed using the non-central F
distribution with non-centrality parameter
\deqn{\lambda = n \cdot k \cdot f^2}{}
where \eqn{n}{} is the per-group sample size, \eqn{k}{} is the number of
groups, and \eqn{f}{} is Cohen's f effect size.
\end{Details}
%
\begin{References}
\begin{enumerate}

\item{} Stroup, W. W., Ptukhina, M., and Garai, S. (2024).
\emph{Generalized Linear Mixed Models: Modern Concepts,
Methods and Applications}. CRC Press.
\item{} Cohen, J. (1988).
\emph{Statistical Power Analysis for the Behavioral Sciences},
2nd ed. Lawrence Erlbaum Associates.

\end{enumerate}

\end{References}
%
\begin{SeeAlso}
\code{\LinkA{Exam21.1}{Exam21.1}}
\end{SeeAlso}
%
\begin{Examples}
\begin{ExampleCode}
data(DataSet21.1)
str(DataSet21.1)
\end{ExampleCode}
\end{Examples}
\HeaderA{DataSet3.1}{Data for Example 3.1 and Example 3.2 from Generalized Linear Mixed Models: Modern Concepts, Methods and Applications by Stroup, Ptukhina, and Garai (2024, 2nd ed.)}{DataSet3.1}
\keyword{datasets}{DataSet3.1}
%
\begin{Description}
DataSet3.1 is used for linear and generalized linear models
\end{Description}
%
\begin{Usage}
\begin{verbatim}
data(DataSet3.1)
\end{verbatim}
\end{Usage}
%
\begin{Format}
A \code{data.frame} with 20 rows and 5 variables.
\end{Format}
%
\begin{Details}
\begin{itemize}

\item{} trt two treatment 0 and 1
\item{} rep unit of observation or observation ID
\item{} Y  is continuous \& may be assumed Gaussian
\item{} N  is the number of obs
\item{} F  is the number of "successes"(N and F specify a binomial response)

\end{itemize}

\end{Details}
%
\begin{Author}
\begin{enumerate}

\item{}  Muhammad Yaseen (\email{myaseen208@gmail.com})
\item{} Adeela Munawar (\email{adeela.uaf@gmail.com})

\end{enumerate}

\end{Author}
%
\begin{References}
\begin{enumerate}

\item{} Stroup, W. W., Ptukhina, M., and Garai, S. (2024).\emph{Generalized linear mixed models: modern concepts, methods and applications}.
CRC Press.

\end{enumerate}

\end{References}
%
\begin{SeeAlso}
\code{\LinkA{Exam3.2}{Exam3.2}}
\end{SeeAlso}
%
\begin{Examples}
\begin{ExampleCode}
data(DataSet3.1)
\end{ExampleCode}
\end{Examples}
\HeaderA{DataSet3.2}{DataSet3.2 for Example 3.3, Example 3.4, Example3.6, Example3.8 and Example 3.9 from Generalized Linear Mixed Models: Modern Concepts, Methods and Applications by Stroup, Ptukhina, and Garai (2024, 2nd ed.)}{DataSet3.2}
\keyword{datasets}{DataSet3.2}
%
\begin{Description}
DataSet3.2 Multi-Location, 4 Treatment Randomized Block
\end{Description}
%
\begin{Usage}
\begin{verbatim}
data(DataSet3.2)
\end{verbatim}
\end{Usage}
%
\begin{Format}
A \code{data.frame} with 32 rows and 10 variables.
\end{Format}
%
\begin{Details}
\begin{itemize}

\item{} trt              four treatments, coded 0, 1, 2, and 3
\item{} loc              eight locations used as blocks
\item{} Y                is Gaussian response variable
\item{} Nbin             subjects at each Loc x Trt for binomial response
\item{} S1 and S2        are two binomial response variables
\item{} count1 and count 2 used later
\item{} A and B          factorial treatment components with levels 0 and 1;
the mapping follows the Chapter 3, Section 3.3.5
\code{a*b} least-squares means table

\end{itemize}

\end{Details}
%
\begin{Author}
\begin{enumerate}

\item{}  Muhammad Yaseen (\email{myaseen208@gmail.com})
\item{} Adeela Munawar (\email{adeela.uaf@gmail.com})

\end{enumerate}

\end{Author}
%
\begin{References}
\begin{enumerate}

\item{} Stroup, W. W., Ptukhina, M., and Garai, S. (2024).\emph{Generalized linear mixed models: modern concepts, methods and applications}.
CRC Press.

\end{enumerate}

\end{References}
%
\begin{SeeAlso}
\code{\LinkA{Exam3.3}{Exam3.3}}
\code{\LinkA{Exam3.9}{Exam3.9}}
\end{SeeAlso}
%
\begin{Examples}
\begin{ExampleCode}
data(DataSet3.2)
\end{ExampleCode}
\end{Examples}
\HeaderA{DataSet3.3}{Data for Example3.7 from Generalized Linear Mixed Models: Modern Concepts, Methods and Applications by Stroup, Ptukhina, and Garai (2024, 2nd ed.)}{DataSet3.3}
\keyword{datasets}{DataSet3.3}
%
\begin{Description}
DataSet3.3 is used to see types of model effects by plotting regression data
\end{Description}
%
\begin{Usage}
\begin{verbatim}
data(DataSet3.3)
\end{verbatim}
\end{Usage}
%
\begin{Format}
A \code{data.frame} with 36 rows and 6 variables.
\end{Format}
%
\begin{Details}
\begin{itemize}

\item{} X     Each batch observed at several times:0,3,6,12,24,36,48 months
\item{} Y     continuous variable observed at each level of X
\item{} Fav   number of successes
\item{} N     independent Bernoulli trials
\item{} Batch Batches as 1, 2, 3, 4
\item{} Count binomial response variable

\end{itemize}

\end{Details}
%
\begin{Author}
\begin{enumerate}

\item{}  Muhammad Yaseen (\email{myaseen208@gmail.com})
\item{} Adeela Munawar (\email{adeela.uaf@gmail.com})

\end{enumerate}

\end{Author}
%
\begin{References}
\begin{enumerate}

\item{} Stroup, W. W., Ptukhina, M., and Garai, S. (2024).\emph{Generalized linear mixed models: modern concepts, methods and applications}.
CRC Press.

\end{enumerate}

\end{References}
%
\begin{Examples}
\begin{ExampleCode}
data(DataSet3.3)
\end{ExampleCode}
\end{Examples}
\HeaderA{DataSet8.1}{Data for Example 8.1 from Generalized Linear Mixed Models: Modern Concepts, Methods and Applications by Stroup, Ptukhina, and Garai (2024, 2nd ed.)}{DataSet8.1}
\keyword{datasets}{DataSet8.1}
%
\begin{Description}
Data for Example 8.1 from Generalized Linear Mixed Models: Modern Concepts, Methods and Applications by Stroup, Ptukhina, and Garai (2024, 2nd ed.)
\end{Description}
%
\begin{Usage}
\begin{verbatim}
data(DataSet8.1)
\end{verbatim}
\end{Usage}
%
\begin{Author}
\begin{enumerate}

\item{}  Muhammad Yaseen (\email{myaseen208@gmail.com})
\item{} Adeela Munawar (\email{adeela.uaf@gmail.com})

\end{enumerate}

\end{Author}
%
\begin{References}
\begin{enumerate}

\item{} Stroup, W. W., Ptukhina, M., and Garai, S. (2024).\emph{Generalized linear mixed models: modern concepts, methods and applications}.
CRC Press.

\end{enumerate}

\end{References}
%
\begin{SeeAlso}
\code{\LinkA{Exam9.1}{Exam9.1}}
\end{SeeAlso}
%
\begin{Examples}
\begin{ExampleCode}
data(DataSet8.1)
\end{ExampleCode}
\end{Examples}
\HeaderA{DataSet8.2}{Data for Example 8.2 from Generalized Linear Mixed Models: Modern Concepts, Methods and Applications by Stroup, Ptukhina, and Garai (2024, 2nd ed.)}{DataSet8.2}
\keyword{datasets}{DataSet8.2}
%
\begin{Description}
Data for Example 8.2 from Generalized Linear Mixed Models: Modern Concepts, Methods and Applications by Stroup, Ptukhina, and Garai (2024, 2nd ed.)
\end{Description}
%
\begin{Usage}
\begin{verbatim}
data(DataSet8.2)
\end{verbatim}
\end{Usage}
%
\begin{Author}
\begin{enumerate}

\item{}  Muhammad Yaseen (\email{myaseen208@gmail.com})
\item{} Adeela Munawar (\email{adeela.uaf@gmail.com})

\end{enumerate}

\end{Author}
%
\begin{References}
\begin{enumerate}

\item{} Stroup, W. W., Ptukhina, M., and Garai, S. (2024).\emph{Generalized linear mixed models: modern concepts, methods and applications}.
CRC Press.

\end{enumerate}

\end{References}
%
\begin{SeeAlso}
\code{\LinkA{Exam9.2}{Exam9.2}}
\end{SeeAlso}
%
\begin{Examples}
\begin{ExampleCode}
data(DataSet8.2)
\end{ExampleCode}
\end{Examples}
\HeaderA{DataSet8.3}{Data for Example 8.3 from Generalized Linear Mixed Models: Modern Concepts, Methods and Applications by Stroup, Ptukhina, and Garai (2024, 2nd ed.)}{DataSet8.3}
\keyword{datasets}{DataSet8.3}
%
\begin{Description}
Data for Example 8.3 from Generalized Linear Mixed Models: Modern Concepts, Methods and Applications by Stroup, Ptukhina, and Garai (2024, 2nd ed.)
\end{Description}
%
\begin{Usage}
\begin{verbatim}
data(DataSet8.3)
\end{verbatim}
\end{Usage}
%
\begin{Author}
\begin{enumerate}

\item{}  Muhammad Yaseen (\email{myaseen208@gmail.com})
\item{} Adeela Munawar (\email{adeela.uaf@gmail.com})

\end{enumerate}

\end{Author}
%
\begin{References}
\begin{enumerate}

\item{} Stroup, W. W., Ptukhina, M., and Garai, S. (2024).\emph{Generalized linear mixed models: modern concepts, methods and applications}.
CRC Press.

\end{enumerate}

\end{References}
%
\begin{SeeAlso}
\code{\LinkA{Exam9.3}{Exam9.3}}
\end{SeeAlso}
%
\begin{Examples}
\begin{ExampleCode}
data(DataSet8.3)
\end{ExampleCode}
\end{Examples}
\HeaderA{DataSet8.4}{Data for Example 8.4 from Generalized Linear Mixed Models: Modern Concepts, Methods and Applications by Stroup, Ptukhina, and Garai (2024, 2nd ed.)}{DataSet8.4}
\keyword{datasets}{DataSet8.4}
%
\begin{Description}
Data for Example 8.4 from Generalized Linear Mixed Models: Modern Concepts, Methods and Applications by Stroup, Ptukhina, and Garai (2024, 2nd ed.)
\end{Description}
%
\begin{Usage}
\begin{verbatim}
data(DataSet8.4)
\end{verbatim}
\end{Usage}
%
\begin{Author}
\begin{enumerate}

\item{}  Muhammad Yaseen (\email{myaseen208@gmail.com})
\item{} Adeela Munawar (\email{adeela.uaf@gmail.com})

\end{enumerate}

\end{Author}
%
\begin{References}
\begin{enumerate}

\item{} Stroup, W. W., Ptukhina, M., and Garai, S. (2024).\emph{Generalized linear mixed models: modern concepts, methods and applications}.
CRC Press.

\end{enumerate}

\end{References}
%
\begin{Examples}
\begin{ExampleCode}
data(DataSet8.4)
\end{ExampleCode}
\end{Examples}
\HeaderA{DataSet8.5}{Data for Example 8.4 from Generalized Linear Mixed Models: Modern Concepts, Methods and Applications by Stroup, Ptukhina, and Garai (2024, 2nd ed.)}{DataSet8.5}
\keyword{datasets}{DataSet8.5}
%
\begin{Description}
Data for Example 8.4 from Generalized Linear Mixed Models: Modern Concepts, Methods and Applications by Stroup, Ptukhina, and Garai (2024, 2nd ed.)
\end{Description}
%
\begin{Usage}
\begin{verbatim}
data(DataSet8.5)
\end{verbatim}
\end{Usage}
%
\begin{Author}
\begin{enumerate}

\item{}  Muhammad Yaseen (\email{myaseen208@gmail.com})
\item{} Adeela Munawar (\email{adeela.uaf@gmail.com})

\end{enumerate}

\end{Author}
%
\begin{References}
\begin{enumerate}

\item{} Stroup, W. W., Ptukhina, M., and Garai, S. (2024).\emph{Generalized linear mixed models: modern concepts, methods and applications}.
CRC Press.

\end{enumerate}

\end{References}
%
\begin{Examples}
\begin{ExampleCode}
data(DataSet8.5)
\end{ExampleCode}
\end{Examples}
\HeaderA{DataSet8.6}{Data for Example 8.6 from Generalized Linear Mixed Models: Modern Concepts, Methods and Applications by Stroup, Ptukhina, and Garai (2024, 2nd ed.)}{DataSet8.6}
\keyword{datasets}{DataSet8.6}
%
\begin{Description}
Data for Example 8.6 from Generalized Linear Mixed Models: Modern Concepts, Methods and Applications by Stroup, Ptukhina, and Garai (2024, 2nd ed.)
\end{Description}
%
\begin{Usage}
\begin{verbatim}
data(DataSet8.6)
\end{verbatim}
\end{Usage}
%
\begin{Author}
\begin{enumerate}

\item{}  Muhammad Yaseen (\email{myaseen208@gmail.com})
\item{} Adeela Munawar (\email{adeela.uaf@gmail.com})

\end{enumerate}

\end{Author}
%
\begin{References}
\begin{enumerate}

\item{} Stroup, W. W., Ptukhina, M., and Garai, S. (2024).\emph{Generalized linear mixed models: modern concepts, methods and applications}.
CRC Press.

\end{enumerate}

\end{References}
%
\begin{SeeAlso}
\code{\LinkA{Exam8.6}{Exam8.6}}
\end{SeeAlso}
%
\begin{Examples}
\begin{ExampleCode}
data(DataSet8.6)
\end{ExampleCode}
\end{Examples}
\HeaderA{DataSet8.7}{Data for Example 8.7 from Generalized Linear Mixed Models: Modern Concepts, Methods and Applications by Stroup, Ptukhina, and Garai (2024, 2nd ed.)}{DataSet8.7}
\keyword{datasets}{DataSet8.7}
%
\begin{Description}
Data for Example 8.7 from Generalized Linear Mixed Models: Modern Concepts, Methods and Applications by Stroup, Ptukhina, and Garai (2024, 2nd ed.)
\end{Description}
%
\begin{Usage}
\begin{verbatim}
data(DataSet8.7)
\end{verbatim}
\end{Usage}
%
\begin{Author}
\begin{enumerate}

\item{}  Muhammad Yaseen (\email{myaseen208@gmail.com})
\item{} Adeela Munawar (\email{adeela.uaf@gmail.com})

\end{enumerate}

\end{Author}
%
\begin{References}
\begin{enumerate}

\item{} Stroup, W. W., Ptukhina, M., and Garai, S. (2024).\emph{Generalized linear mixed models: modern concepts, methods and applications}.
CRC Press.

\end{enumerate}

\end{References}
%
\begin{Examples}
\begin{ExampleCode}
data(DataSet8.7)
\end{ExampleCode}
\end{Examples}
\HeaderA{DataSet9.1}{Data for Example 9.1 from Generalized Linear Mixed Models: Modern Concepts, Methods and Applications by Stroup, Ptukhina, and Garai (2024, 2nd ed.)}{DataSet9.1}
\keyword{datasets}{DataSet9.1}
%
\begin{Description}
DataSet9.1 is used for Nested factorial structure
\end{Description}
%
\begin{Usage}
\begin{verbatim}
data(DataSet9.1)
\end{verbatim}
\end{Usage}
%
\begin{Format}
A \code{data.frame} with 30 rows and 4 variables.
\end{Format}
%
\begin{Details}
\begin{itemize}

\item{} block  10 blocks
\item{} trt    6 treatments nested within sets
\item{} set    2 sets
\item{} y      is a Gaussian response variable

\end{itemize}

\end{Details}
%
\begin{Author}
Muhammad Yaseen (\email{myaseen208@gmail.com})
Adeela Munawar (\email{adeela.uaf@gmail.com})
\end{Author}
%
\begin{References}
\begin{enumerate}

\item{} Stroup, W. W., Ptukhina, M., and Garai, S. (2024).\emph{Generalized Linear Mixed Models: Modern Concepts, Methods and Applications (2nd ed.)}.
CRC Press.

\end{enumerate}

\end{References}
%
\begin{SeeAlso}
\code{\LinkA{Exam9.1}{Exam9.1}}
\end{SeeAlso}
%
\begin{Examples}
\begin{ExampleCode}
data(DataSet9.1)
\end{ExampleCode}
\end{Examples}
\HeaderA{DataSet9.2}{Data for Example 9.2 from Generalized Linear Mixed Models: Modern Concepts, Methods and Applications by Stroup, Ptukhina, and Garai (2024, 2nd ed.)}{DataSet9.2}
\keyword{datasets}{DataSet9.2}
%
\begin{Description}
DataSet9.2 is used for Incomplete strip-plot ( 3 cross 3 factorial).
\end{Description}
%
\begin{Usage}
\begin{verbatim}
data(DataSet9.2)
\end{verbatim}
\end{Usage}
%
\begin{Format}
A \code{data.frame} with 36 rows and 6 variables.
\end{Format}
%
\begin{Details}
\begin{itemize}

\item{} block  9 blocks each consisting of 2 rows and 2 columns
\item{} a      is a factor with 3 levels assigned at random to rows
\item{} b      is a factor with 3 levels assigned at random to columns
\item{} y      is a Gaussian response variable

\end{itemize}

\end{Details}
%
\begin{Author}
Muhammad Yaseen (\email{myaseen208@gmail.com})
Adeela Munawar (\email{adeela.uaf@gmail.com})
\end{Author}
%
\begin{References}
\begin{enumerate}

\item{} Stroup, W. W., Ptukhina, M., and Garai, S. (2024).\emph{Generalized Linear Mixed Models: Modern Concepts, Methods and Applications (2nd ed.)}.
CRC Press.

\end{enumerate}

\end{References}
%
\begin{SeeAlso}
\code{\LinkA{Exam9.2}{Exam9.2}}
\end{SeeAlso}
%
\begin{Examples}
\begin{ExampleCode}
data(DataSet9.2)
\end{ExampleCode}
\end{Examples}
\HeaderA{DataSet9.3}{Data for Example 9.3 from Generalized Linear Mixed Models: Modern Concepts, Methods and Applications by Stroup, Ptukhina, and Garai (2024, 2nd ed.)}{DataSet9.3}
\keyword{datasets}{DataSet9.3}
%
\begin{Description}
DataSet9.3 is used for Response surface design with incomplete blocking
\end{Description}
%
\begin{Usage}
\begin{verbatim}
data(DataSet9.3)
\end{verbatim}
\end{Usage}
%
\begin{Format}
A \code{data.frame} with 28 rows and 4 variables.
\end{Format}
%
\begin{Details}
\begin{itemize}

\item{} block  with 7 blocks
\item{} a      is a factor with 3 levels 0,-1 and 1
\item{} b      is a factor with 3 levels 0,-1 and 1
\item{} c      is a factor with 3 levels 0,-1 and 1
\item{} y      is a Gaussian response variable

\end{itemize}

\end{Details}
%
\begin{Author}
Muhammad Yaseen (\email{myaseen208@gmail.com})
Adeela Munawar (\email{adeela.uaf@gmail.com})
\end{Author}
%
\begin{References}
\begin{enumerate}

\item{} Stroup, W. W., Ptukhina, M., and Garai, S. (2024).\emph{Generalized Linear Mixed Models: Modern Concepts, Methods and Applications (2nd ed.)}.
CRC Press.

\end{enumerate}

\end{References}
%
\begin{SeeAlso}
\code{\LinkA{Exam9.3}{Exam9.3}}
\end{SeeAlso}
%
\begin{Examples}
\begin{ExampleCode}
data(DataSet9.3)
\end{ExampleCode}
\end{Examples}
\HeaderA{DataSet9.4}{Data for Example 9.4 from Generalized Linear Mixed Models: Modern Concepts, Methods and Applications by Stroup, Ptukhina, and Garai (2024, 2nd ed.)}{DataSet9.4}
\keyword{datasets}{DataSet9.4}
%
\begin{Description}
DataSet9.4 is used for Multifactor treatment and Multilevel design structures
\end{Description}
%
\begin{Usage}
\begin{verbatim}
data(DataSet9.4)
\end{verbatim}
\end{Usage}
%
\begin{Format}
A \code{data.frame} with 36 rows and 6 variables.
\end{Format}
%
\begin{Details}
\begin{itemize}

\item{} block  9 blocks each consisting of 2 rows and 2 columns
\item{} a      is a factor with 3 levels assigned at random to rows
\item{} b      is a factor with 3 levels assigned at random to columns
\item{} y      is a Gaussian response variable

\end{itemize}

\end{Details}
%
\begin{Author}
Muhammad Yaseen (\email{myaseen208@gmail.com})
Adeela Munawar (\email{adeela.uaf@gmail.com})
\end{Author}
%
\begin{References}
\begin{enumerate}

\item{} Stroup, W. W., Ptukhina, M., and Garai, S. (2024).\emph{Generalized Linear Mixed Models: Modern Concepts, Methods and Applications (2nd ed.)}.
CRC Press.

\end{enumerate}

\end{References}
%
\begin{SeeAlso}
\code{\LinkA{Exam9.4}{Exam9.4}}
\end{SeeAlso}
%
\begin{Examples}
\begin{ExampleCode}
data(DataSet9.4)
\end{ExampleCode}
\end{Examples}
\HeaderA{Exam1.1}{Example 1.1 — Probability Distribution, LM, and GLM}{Exam1.1}
%
\begin{Description}
Example 1.1 (Stroup et al., 2024) illustrates the difference between fitting
a Gaussian linear model and a logistic GLM to binary proportion data.
Using Table 1.1 dose–response data, it shows that the logistic GLM provides
better fit to bounded proportions than a linear model that can exceed [0,1].
\end{Description}
%
\begin{Author}
\begin{enumerate}

\item{} Muhammad Yaseen (\email{myaseen208@gmail.com})
\item{} Adeela Munawar (\email{adeela.uaf@gmail.com})

\end{enumerate}

\end{Author}
%
\begin{References}
\begin{enumerate}

\item{} Stroup, W. W., Ptukhina, M., and Garai, S. (2024).
\emph{Generalized Linear Mixed Models: Modern Concepts, Methods and Applications (2nd ed.)}.
CRC Press.

\end{enumerate}

\end{References}
%
\begin{SeeAlso}
\LinkA{Table1.1}{Table1.1}
\end{SeeAlso}
%
\begin{Examples}
\begin{ExampleCode}
data(Table1.1)

## Linear Model (Section 1.2.3)
Exam1.1.lm1 <- stats::lm(formula = y / Nx ~ x, data = Table1.1)
summary(Exam1.1.lm1)
if (requireNamespace("parameters", quietly = TRUE)) {
  parameters::model_parameters(Exam1.1.lm1)
}

## GLM with logit link (binomial family, Section 1.2.4)
Exam1.1.glm1 <- stats::glm(
  formula = cbind(y, Nx - y) ~ x,
  family  = stats::binomial(link = "logit"),
  data    = Table1.1
)
summary(Exam1.1.glm1)
if (requireNamespace("parameters", quietly = TRUE)) {
  parameters::model_parameters(Exam1.1.glm1)
}

## Figure 1.1 — LM vs GLM fitted proportions
plot_df <- rbind(
  data.frame(x = Table1.1$x,
             p = Table1.1$y / Table1.1$Nx,
             Model = "Observed"),
  data.frame(x = Table1.1$x,
             p = Exam1.1.lm1$fitted.values,
             Model = "LM"),
  data.frame(x = Table1.1$x,
             p = Exam1.1.glm1$fitted.values,
             Model = "GLM (logit)")
)

ggplot2::ggplot(
    plot_df,
    ggplot2::aes(x = x, y = p, colour = Model, shape = Model)
  ) +
  ggplot2::geom_point(size = 2.5) +
  ggplot2::geom_line(
    data = subset(plot_df, Model != "Observed")
  ) +
  ggplot2::scale_colour_manual(
    values = c("Observed"    = "black",
               "LM"          = "#377EB8",
               "GLM (logit)" = "#E41A1C")
  ) +
  ggplot2::labs(
    title  = "Example 1.1: Linear Model vs Logistic GLM",
    x      = "Dose (x)",
    y      = "Proportion (p)",
    colour = "Model",
    shape  = "Model"
  ) +
  ggplot2::theme_bw()

## Correlation of fitted values with observed
cat("LM  correlation:", round(stats::cor(Table1.1$y / Table1.1$Nx,
                                        Exam1.1.lm1$fitted.values), 4), "\n")
cat("GLM correlation:", round(stats::cor(Table1.1$y / Table1.1$Nx,
                                        Exam1.1.glm1$fitted.values), 4), "\n")
\end{ExampleCode}
\end{Examples}
\HeaderA{Exam1.2}{Example1.2 from Generalized Linear Mixed Models: Modern Concepts, Methods and Applications by Stroup, Ptukhina, and Garai (2024, 2nd ed.)}{Exam1.2}
%
\begin{Description}
Exam1.2 is used to see types of model effects by plotting regression data
\end{Description}
%
\begin{Author}
\begin{enumerate}

\item{}  Muhammad Yaseen (\email{myaseen208@gmail.com})
\item{} Adeela Munawar (\email{adeela.uaf@gmail.com})

\end{enumerate}

\end{Author}
%
\begin{References}
\begin{enumerate}

\item{} Stroup, W. W., Ptukhina, M., and Garai, S. (2024).
\emph{Generalized Linear Mixed Models: Modern Concepts, Methods and Applications (2nd ed.)}.
CRC Press.

\end{enumerate}

\end{References}
%
\begin{SeeAlso}
\code{\LinkA{Table1.2}{Table1.2}}
\end{SeeAlso}
%
\begin{Examples}
\begin{ExampleCode}

#-------------------------------------------------------------
## Plot of multi-batch regression data discussed in Article 1.3
#-------------------------------------------------------------
data(Table1.2)

Table1.2$Batch <- factor(x  = Table1.2$Batch)

Plot  <-
 ggplot2::ggplot(
   data = Table1.2,
   mapping = ggplot2::aes(y = Y, x = X, colour = Batch, shape = Batch)
 ) +
 ggplot2::geom_point() +
 ggplot2::geom_smooth(method = "lm", fill =  NA) +
 ggplot2::labs(title   = "Plot of Multi Batch Regression data") +
 ggplot2::theme_bw()
Plot
\end{ExampleCode}
\end{Examples}
\HeaderA{Exam10.1}{Example 10.1 from Generalized Linear Mixed Models: Modern Concepts, Methods and Applications by Stroup, Ptukhina, and Garai (2024, 2nd ed.)}{Exam10.1}
%
\begin{Description}
Exam10.1 One-way random effects model and BLUP estimation
\end{Description}
%
\begin{Author}
\begin{enumerate}

\item{}  Muhammad Yaseen (\email{myaseen208@gmail.com})
\item{} Adeela Munawar (\email{adeela.uaf@gmail.com})

\end{enumerate}

\end{Author}
%
\begin{References}
\begin{enumerate}

\item{} Stroup, W. W., Ptukhina, M., and Garai, S. (2024).
\emph{Generalized Linear Mixed Models: Modern Concepts, Methods and Applications (2nd ed.)}.
CRC Press.

\end{enumerate}

\end{References}
%
\begin{SeeAlso}
\code{\LinkA{DataSet10.1}{DataSet10.1}}
\end{SeeAlso}
%
\begin{Examples}
\begin{ExampleCode}

data(DataSet10.1)
DataSet10.1$a <- factor(x = DataSet10.1$a)

## Random effects model
Exam10.1lmer <- lme4::lmer(y ~ 1 + (1|a), data = DataSet10.1)
summary(Exam10.1lmer)

## Fixed effects model (for comparison)
Exam10.1lm <- stats::lm(y ~ a, data = DataSet10.1)
summary(Exam10.1lm)

## Overall mean — broad inference (random effects model)
emmeans::emmeans(Exam10.1lmer, specs = ~1)

## Overall mean — narrow inference (fixed effects model)
emmeans::emmeans(Exam10.1lm, specs = ~1)

## BLUP Estimates and standard errors for each level of a
## SE from conditional variance: sqrt(postVar) matches book Table 10.1 SE = 0.82453
blup_coef <- unlist(lme4::ranef(Exam10.1lmer)$a)
rv_a      <- lme4::ranef(Exam10.1lmer, condVar = TRUE)
BLUPa_SE  <- sqrt(attr(rv_a$a, "postVar")[1, 1, ])
BLUPa     <- mean(DataSet10.1$y) + blup_coef
print(data.frame(level = rownames(rv_a$a), BLUP = BLUPa, SE = BLUPa_SE))

\end{ExampleCode}
\end{Examples}
\HeaderA{Exam10.2}{Example 10.2 from Generalized Linear Mixed Models: Modern Concepts, Methods and Applications by Stroup, Ptukhina, and Garai (2024, 2nd ed.)}{Exam10.2}
%
\begin{Description}
Exam10.2 Two-way nested random effects model and BLUP estimation
\end{Description}
%
\begin{Author}
\begin{enumerate}

\item{}  Muhammad Yaseen (\email{myaseen208@gmail.com})
\item{} Adeela Munawar (\email{adeela.uaf@gmail.com})

\end{enumerate}

\end{Author}
%
\begin{References}
\begin{enumerate}

\item{} Stroup, W. W., Ptukhina, M., and Garai, S. (2024).
\emph{Generalized Linear Mixed Models: Modern Concepts, Methods and Applications (2nd ed.)}.
CRC Press.

\end{enumerate}

\end{References}
%
\begin{SeeAlso}
\code{\LinkA{DataSet10.2}{DataSet10.2}}
\end{SeeAlso}
%
\begin{Examples}
\begin{ExampleCode}

data(DataSet10.2)
DataSet10.2$a <- factor(x = DataSet10.2$a)
DataSet10.2$b <- factor(x = DataSet10.2$b)

## Random effects nested model (b nested within a)
Exam10.2lmer <- lmerTest::lmer(y ~ 1 + (1 | a / b), data = DataSet10.2)
summary(Exam10.2lmer)

## Fixed effects nested model (for comparison)
Exam10.2lm <- stats::lm(y ~ a + b %in% a, data = DataSet10.2)
summary(Exam10.2lm)

## Overall mean — narrow inference
emmeans::emmeans(Exam10.2lm, specs = ~1)

## BLUP Estimates for each level of a
blup_coef <- unlist(lme4::ranef(Exam10.2lmer)$a)
BLUPa <- sapply(seq_along(blup_coef), \(i) mean(DataSet10.2$y) + blup_coef[i])
BLUPa

\end{ExampleCode}
\end{Examples}
\HeaderA{Exam10.4}{Example 10.4 from Generalized Linear Mixed Models: Modern Concepts, Methods and Applications by Stroup, Ptukhina, and Garai (2024, 2nd ed.)}{Exam10.4}
%
\begin{Description}
Exam10.4 Relationship between BLUP and Fixed Effect Estimators
\end{Description}
%
\begin{Author}
\begin{enumerate}

\item{}  Muhammad Yaseen (\email{myaseen208@gmail.com})
\item{} Adeela Munawar (\email{adeela.uaf@gmail.com})

\end{enumerate}

\end{Author}
%
\begin{References}
\begin{enumerate}

\item{} Stroup, W. W., Ptukhina, M., and Garai, S. (2024).
\emph{Generalized Linear Mixed Models: Modern Concepts, Methods and Applications (2nd ed.)}.
CRC Press.

\end{enumerate}

\end{References}
%
\begin{SeeAlso}
\code{\LinkA{DataSet10.4}{DataSet10.4}}
\end{SeeAlso}
%
\begin{Examples}
\begin{ExampleCode}

data(DataSet10.4)
DataSet10.4$a <- factor(x = DataSet10.4$a)
DataSet10.4$b <- factor(x = DataSet10.4$b)

Exam10.4lmer <- lme4::lmer(y ~ a + (1|b) + (1|a:b), data = DataSet10.4)
summary(Exam10.4lmer)

## BLUP vs fixed effects comparison
emmeans::emmeans(Exam10.4lmer, specs = ~a)

\end{ExampleCode}
\end{Examples}
\HeaderA{Exam11.1}{Example 11.1 - CRD Count Data and Poisson-Normal GLMM}{Exam11.1}
%
\begin{Description}
Reproduces the Chapter 11, Section 11.1.2 workflow comparing an ANOVA on
counts, an ANOVA on log counts, a Poisson GLM, and a Poisson-normal GLMM
with an observation-level random effect.
\end{Description}
%
\begin{Details}
The Poisson-normal GLMM is
\deqn{\log(\lambda_{ij}) = \eta + \tau_i + u_{ij}, \quad
      u_{ij} \sim N(0, \sigma_U^2).}{}
\end{Details}
%
\begin{Note}
The count and log-count ANOVA outputs reproduce the printed book values.
The Poisson-normal GLMM is fitted with \pkg{lme4}; small differences from
SAS PROC GLIMMIX are expected because the estimation algorithms differ.
\end{Note}
%
\begin{References}
Stroup, W. W., Ptukhina, M., and Garai, S. (2024).
\emph{Generalized Linear Mixed Models: Modern Concepts, Methods and
Applications}, 2nd ed. CRC Press.
\end{References}
%
\begin{SeeAlso}
\code{\LinkA{DataSet11.1}{DataSet11.1}}
\end{SeeAlso}
%
\begin{Examples}
\begin{ExampleCode}
data(DataSet11.1)

fit_count <- stats::lm(count ~ trt, data = DataSet11.1)
stats::anova(fit_count)

DataSet11.1$log_count <- log(DataSet11.1$count)
fit_log <- stats::lm(log_count ~ trt, data = DataSet11.1)
stats::anova(fit_log)
emmeans::emmeans(fit_log, specs = ~ trt)

fit_glm <- stats::glm(
  count ~ trt,
  family = stats::poisson(link = "log"),
  data = DataSet11.1
)
emmeans::emmeans(fit_glm, specs = ~ trt, type = "response")

if (requireNamespace("lme4", quietly = TRUE)) {
  fit_glmm <- lme4::glmer(
    count ~ trt + (1 | trt:unit),
    family = stats::poisson(link = "log"),
    data = DataSet11.1,
    nAGQ = 0,
    control = lme4::glmerControl(optimizer = "bobyqa")
  )
  summary(fit_glmm)
  lme4::VarCorr(fit_glmm)
  emmeans::emmeans(fit_glmm, specs = ~ trt, type = "response")

  if (requireNamespace("DHARMa", quietly = TRUE)) {
    sim_res <- DHARMa::simulateResiduals(fit_glmm, plot = FALSE)
    DHARMa::testDispersion(sim_res)
  }

  if (requireNamespace("report", quietly = TRUE)) {
    report::report(fit_glmm)
  }
}
\end{ExampleCode}
\end{Examples}
\HeaderA{Exam11.3}{Example 11.3 - Blocked Count Data and Overdispersion}{Exam11.3}
%
\begin{Description}
Fits the Chapter 11 blocked-design count-data models from Section 11.4:
a naive Poisson GLMM, a Poisson-normal GLMM with a block-by-treatment unit
effect, and a negative-binomial GLMM.
\end{Description}
%
\begin{Details}
The naive Poisson model is
\deqn{\log(\lambda_{ij}) = \eta + \tau_i + r_j,\quad
      r_j \sim N(0,\sigma_R^2).}{}
The Poisson-normal model adds a unit term
\deqn{\log(\lambda_{ij}) = \eta + \tau_i + r_j + (rt)_{ij}.}{}
The negative-binomial model omits the unit term and uses the distributional
scale parameter to account for unit-level overdispersion.
\end{Details}
%
\begin{Note}
\pkg{glmmTMB} uses the \code{nbinom2} parameterization
\eqn{Var(Y) = \mu + \mu^2/\theta}{}. The book's negative-binomial scale
\eqn{\phi}{} corresponds approximately to \eqn{1/\theta}{}.
\end{Note}
%
\begin{References}
Stroup, W. W., Ptukhina, M., and Garai, S. (2024).
\emph{Generalized Linear Mixed Models: Modern Concepts, Methods and
Applications}, 2nd ed. CRC Press.
\end{References}
%
\begin{SeeAlso}
\code{\LinkA{DataSet11.3}{DataSet11.3}}
\end{SeeAlso}
%
\begin{Examples}
\begin{ExampleCode}
data(DataSet11.3)

if (requireNamespace("lme4", quietly = TRUE)) {
  fit_naive <- lme4::glmer(
    count ~ trt + (1 | block),
    family = stats::poisson(link = "log"),
    data = DataSet11.3,
    nAGQ = 1,
    control = lme4::glmerControl(optimizer = "bobyqa")
  )
  pearson_df <- sum(stats::residuals(fit_naive, type = "pearson")^2) /
    stats::df.residual(fit_naive)
  pearson_df
  emmeans::emmeans(fit_naive, specs = ~ trt, type = "response")

  fit_pois_normal <- lme4::glmer(
    count ~ trt + (1 | block) + (1 | block:trt),
    family = stats::poisson(link = "log"),
    data = DataSet11.3,
    nAGQ = 0,
    control = lme4::glmerControl(optimizer = "bobyqa")
  )
  lme4::VarCorr(fit_pois_normal)
  emmeans::emmeans(fit_pois_normal, specs = ~ trt, type = "response")
}

if (requireNamespace("glmmTMB", quietly = TRUE)) {
  fit_nb <- glmmTMB::glmmTMB(
    count ~ trt + (1 | block),
    family = glmmTMB::nbinom2(link = "log"),
    data = DataSet11.3
  )
  summary(fit_nb)
  emmeans::emmeans(fit_nb, specs = ~ trt, type = "response")

  if (requireNamespace("DHARMa", quietly = TRUE)) {
    sim_res <- DHARMa::simulateResiduals(fit_nb, plot = FALSE)
    DHARMa::testDispersion(sim_res)
  }

  if (requireNamespace("report", quietly = TRUE)) {
    report::report(fit_nb)
  }
}
\end{ExampleCode}
\end{Examples}
\HeaderA{Exam11.4}{Example 11.4 - Split-Plot Count Data}{Exam11.4}
%
\begin{Description}
Fits the Chapter 11, Section 11.5 split-plot count-data model for a
7 by 4 factorial treatment structure in four blocks.
\end{Description}
%
\begin{Details}
The negative-binomial GLMM for the conditional inference target is
\deqn{\log(\lambda_{ijk}) = \eta + \alpha_i + \beta_j +
      (\alpha\beta)_{ij} + r_k + (ra)_{ik}.}{}
\end{Details}
%
\begin{Note}
The package dataset is synthetic because the original \code{sp\_counts}
SAS data were not available. Use this example as an executable R analogue
of the Chapter 11 workflow rather than as a numerical reproduction of all
book tables.
\end{Note}
%
\begin{References}
Stroup, W. W., Ptukhina, M., and Garai, S. (2024).
\emph{Generalized Linear Mixed Models: Modern Concepts, Methods and
Applications}, 2nd ed. CRC Press.
\end{References}
%
\begin{SeeAlso}
\code{\LinkA{DataSet11.4}{DataSet11.4}}
\end{SeeAlso}
%
\begin{Examples}
\begin{ExampleCode}
data(DataSet11.4)

if (requireNamespace("glmmTMB", quietly = TRUE)) {
  fit_nb <- glmmTMB::glmmTMB(
    count ~ a * b + (1 | block) + (1 | block:a),
    family = glmmTMB::nbinom2(link = "log"),
    data = DataSet11.4
  )
  summary(fit_nb)
  emmeans::emmeans(fit_nb, specs = ~ a * b, type = "response")

  if (requireNamespace("DHARMa", quietly = TRUE)) {
    sim_res <- DHARMa::simulateResiduals(fit_nb, plot = FALSE)
    DHARMa::testDispersion(sim_res)
  }

  if (requireNamespace("report", quietly = TRUE)) {
    report::report(fit_nb)
  }
}
\end{ExampleCode}
\end{Examples}
\HeaderA{Exam12.1}{Example 12.1 — Continuous Proportions: Beta GLMM}{Exam12.1}
%
\begin{Description}
Analysis of continuous proportion data from a dose-response study using
a beta GLMM with a run-level random effect.  Implements the analysis of
Section 12.6.2 of Stroup et al. (2024), Data Set 12.7.  Two treatments
are compared across six dose levels (0–5), with 12 runs per treatment.
Demonstrates lack-of-fit testing, regression slope comparisons, and
treatment contrasts at specific dose levels.
\end{Description}
%
\begin{Details}
The beta GLMM for continuous proportions is:
\deqn{y_{ijk} | r(\tau)_{ik} \sim \text{Beta}(\mu_{ijk}, \phi)}{}
\deqn{\text{logit}(\mu_{ijk}) = \beta_{0i} + \beta_{1i} D_j + r(\tau)_{ik}}{}
where \eqn{\beta_{0i}}{} is the treatment intercept, \eqn{\beta_{1i}}{} is the
dose-specific linear slope, \eqn{D_j}{} is the dose level (0–5), and
\eqn{r(\tau)_{ik} \sim \mathcal{N}(0,\sigma^2_r)}{} is the run random
effect nested within treatment.

The key hypotheses are:
\begin{itemize}

\item{} \eqn{H_0: \beta_{0,0} = \beta_{0,1}}{} (equal intercepts)
\item{} \eqn{H_0: \beta_{1,0} = \beta_{1,1}}{} (equal slopes)

\end{itemize}


Published results (2nd ed. pp.406-407): equal intercepts F(1,22)=0.17,
p=0.6837; equal slopes F(1,118)=10.07, p=0.0019.
\end{Details}
%
\begin{Note}
\strong{Dataset note:} \code{DataSet12.1} is reconstructed from the
published regression parameters (p.406) with run RE drawn from
\eqn{N(0,0.45^2)}{} and Beta precision \eqn{\phi=10}{} (\code{set.seed(2024)}).
The actual data appear in the SAS Data and Program Library as Data Set 12.7
(SAS name \code{rates}).  The contrast conclusions
(non-significant equal intercepts, significant unequal slopes) match the
book; individual fixed-effect estimates approximate the published values.
\end{Note}
%
\begin{References}
\begin{enumerate}

\item{} Stroup, W. W., Ptukhina, M., and Garai, S. (2024).
\emph{Generalized Linear Mixed Models: Modern Concepts, Methods and
Applications}. CRC Press. Section 12.6.2, pp.404–407.
\item{} Ferrari, S., \& Cribari-Neto, F. (2004).
Beta regression for modelling rates and proportions.
\emph{Journal of Applied Statistics}, 31(7), 799–815.

\end{enumerate}

\end{References}
%
\begin{SeeAlso}
\code{\LinkA{DataSet12.1}{DataSet12.1}}
\end{SeeAlso}
%
\begin{Examples}
\begin{ExampleCode}
data(DataSet12.1)

## -------------------------------------------------------
## 1. Exploratory: mean proportion by treatment and dose
## -------------------------------------------------------
with(DataSet12.1, tapply(proportion, list(trt, dose), mean))

## -------------------------------------------------------
## 2. Beta GLMM with run(trt) random effect via glmmTMB
## -------------------------------------------------------
if (requireNamespace("glmmTMB", quietly = TRUE)) {

  ## Step 1: lack-of-fit test (add trt×dose interaction)
  fit_full <- glmmTMB::glmmTMB(
    proportion ~ trt * dose + (1 | run),
    family = glmmTMB::beta_family(link = "logit"),
    data   = DataSet12.1
  )
  summary(fit_full)

  ## Step 2: separate linear regressions per treatment
  fit_sep <- glmmTMB::glmmTMB(
    proportion ~ trt / dose - 1 + (1 | run),
    family = glmmTMB::beta_family(link = "logit"),
    data   = DataSet12.1
  )
  summary(fit_sep)

  ## -------------------------------------------------------
  ## 3. Contrast: equal intercepts and equal slopes
  ## -------------------------------------------------------
  emm_int <- emmeans::emmeans(fit_sep, ~ trt, at = list(dose = 0))
  emmeans::contrast(emm_int, method = "pairwise")

  emm_slope <- emmeans::emtrends(fit_sep, ~ trt, var = "dose")
  emmeans::contrast(emm_slope, method = "pairwise")

  ## -------------------------------------------------------
  ## 4. Predicted proportions at dose=0 and dose=5
  ## -------------------------------------------------------
  emm_d0 <- emmeans::emmeans(fit_sep, ~ trt, at = list(dose = 0),
                              type = "response")
  emm_d5 <- emmeans::emmeans(fit_sep, ~ trt, at = list(dose = 5),
                              type = "response")
  print(emm_d0)
  print(emm_d5)
  emmeans::contrast(emm_d5, method = "pairwise")
}
\end{ExampleCode}
\end{Examples}
\HeaderA{Exam12.2}{Example 12.2 — Binomial GLMM: Nested Factorial Design}{Exam12.2}
%
\begin{Description}
Analysis of binomial count data from a nested factorial design (10 blocks,
6 treatments in 2 sets, 3 treatments per set) using a GLMM with binomial
response and logit link.  Implements Section 12.3.2 of Stroup et al. (2024),
Data Set 12.2.  Random block effects are modelled with a random intercept
plus a random set-by-block slope to allow varying set effects across blocks.
\end{Description}
%
\begin{Details}
The nested factorial GLMM is:
\deqn{\text{logit}(\pi_{ijk}) = \eta + \alpha_i + \beta_{j(i)} + r_{k(i)}}{}
where \eqn{\alpha_i}{} denotes the set (A) effect, \eqn{\beta_{j(i)}}{}
denotes the treatment B nested within set, and
\eqn{r_{k(i)} \sim \mathcal{N}(0,\sigma^2_A)}{} is the block-within-A
random effect (whole-plot error for the nested factorial).  Blocks are
nested within A: blocks 1–5 receive setA, blocks 6–10 receive setB.

Published results (2nd ed. p.382):
\eqn{\hat\sigma^2_A = 0.4784}{} (SE=0.2664);
A: F(1,8)=0.46, p=0.5187; B(A): F(4,16)=2.43, p=0.0907.
\end{Details}
%
\begin{Note}
\strong{Dataset note:} \code{DataSet12.2} is reconstructed from the
published design description (5 blocks × setA + 5 blocks × setB, 3
treatments per block, N=20 per cell, 30 obs total) using the published
logit B(A) LSMeans (p.382) and block variance \eqn{\sigma^2_A=0.4784}{}
as generative parameters.  The actual data appear in the SAS Data and
Program Library as Data Set 12.2 (SAS name \code{ch11\_ex\_12C2}).
Numerical output will approximate but not exactly reproduce published
values.
\end{Note}
%
\begin{References}
\begin{enumerate}

\item{} Stroup, W. W., Ptukhina, M., and Garai, S. (2024).
\emph{Generalized Linear Mixed Models: Modern Concepts, Methods and
Applications}. CRC Press. Sections 12.3.2 and 12.4.1, pp.380–393.

\end{enumerate}

\end{References}
%
\begin{SeeAlso}
\code{\LinkA{DataSet12.2}{DataSet12.2}}
\end{SeeAlso}
%
\begin{Examples}
\begin{ExampleCode}
data(DataSet12.2)

## -------------------------------------------------------
## 1. Exploratory: observed proportions by set and treatment
## -------------------------------------------------------
with(DataSet12.2, tapply(f / n, list(a, b), mean))

## -------------------------------------------------------
## 2. Binomial GLMM — nested factorial
## Blocks are nested within A: (1 | block) = whole-plot error
## Matches book random a / subject=Block in GLIMMIX
## -------------------------------------------------------
fit_lme4 <- lme4::glmer(
  cbind(f, n - f) ~ a / b + (1 | block),
  family  = stats::binomial(link = "logit"),
  data    = DataSet12.2,
  control = lme4::glmerControl(optimizer = "bobyqa")
)
summary(fit_lme4)

## -------------------------------------------------------
## 3. Type III tests for A and B(A)
## -------------------------------------------------------
if (requireNamespace("car", quietly = TRUE)) {
  car::Anova(fit_lme4, type = "III")
}

## -------------------------------------------------------
## 4. Least-squares means for B(A) — treatment within set
## -------------------------------------------------------
emm_ba <- emmeans::emmeans(fit_lme4, ~ b | a, type = "response")
print(emm_ba)
emmeans::contrast(emm_ba, method = "pairwise")

## -------------------------------------------------------
## 5. Probit link (Section 12.4.1)
## -------------------------------------------------------
fit_probit <- lme4::glmer(
  cbind(f, n - f) ~ a / b + (1 | block),
  family  = stats::binomial(link = "probit"),
  data    = DataSet12.2,
  control = lme4::glmerControl(optimizer = "bobyqa")
)
summary(fit_probit)

## Compare logit vs probit: F-test for B(A) within A=setA
emm_probit <- emmeans::emmeans(fit_probit, ~ b | a, type = "response")
print(emm_probit)

## -------------------------------------------------------
## 6. Overdispersion check
## -------------------------------------------------------
if (requireNamespace("DHARMa", quietly = TRUE)) {
  sim_r <- DHARMa::simulateResiduals(fit_lme4)
  DHARMa::testDispersion(sim_r)
}
\end{ExampleCode}
\end{Examples}
\HeaderA{Exam14.1}{Example 14.1 — Ordinal GLMM: Proportional Odds and Threshold Models}{Exam14.1}
%
\begin{Description}
Analysis of ordinal frequency data from a nested factorial design (10 blocks,
6 treatments, 3-category response) using proportional odds (cumulative logit)
and threshold (cumulative probit) GLMMs.  Implements Section 14.3 of Stroup
et al. (2024), Data Set 14.1.  Each block-treatment combination contributes a
triplet of frequency counts (slight, modrat, severe).
\end{Description}
%
\begin{Details}
The proportional-odds GLMM has one link function per category boundary:
\deqn{\eta_{c,ij} = \eta_c + \tau_i + b_j, \quad c = 0,1}{}
\deqn{\eta_{0,ij} = \log\!\left(\frac{\pi_{\text{slight},ij}}
      {1-\pi_{\text{slight},ij}}\right), \quad
      \eta_{1,ij} = \log\!\left(\frac{\pi_{\text{slight},ij}+
      \pi_{\text{modrat},ij}}{1-(\pi_{\text{slight},ij}+
      \pi_{\text{modrat},ij})}\right)}{}
where \eqn{\eta_c}{} are the category-boundary intercepts (shared across
treatments in the proportional-odds model), \eqn{\tau_i}{} is the treatment
effect, and \eqn{b_j \sim \mathcal{N}(0,\sigma^2_b)}{} is the block random
effect.

Published results (2nd ed. p.436):
\eqn{\hat\eta_0 = 0.3492}{} (slight boundary), \eqn{\hat\eta_1 = 1.9956}{}
(modrat boundary); TRT effects: trt0=\eqn{-2.83}{}, \ldots, trt5=0 (reference);
F(5,768)=17.67, p<0.0001.
\end{Details}
%
\begin{Note}
\strong{Dataset note:} The first 9 rows of \code{DataSet14.1} are transcribed
exactly from book p.435 (blk=1, trt=0,1,2; slight/modrat/severe counts).
Rows 10–180 are reconstructed from the published cumulative-logit parameters.
The actual data appear in the SAS Data and Program Library as Data Set 14.1
(SAS name \code{univar\_multinom}).  Published F and intercept values will
approximate the book's output but differ due to reconstruction.
\end{Note}
%
\begin{References}
\begin{enumerate}

\item{} Stroup, W. W., Ptukhina, M., and Garai, S. (2024).
\emph{Generalized Linear Mixed Models: Modern Concepts, Methods and
Applications}. CRC Press. Section 14.3, pp.434–438.
\item{} Agresti, A. (2010). \emph{Analysis of Ordinal Categorical Data},
2nd ed. Wiley.

\end{enumerate}

\end{References}
%
\begin{SeeAlso}
\code{\LinkA{DataSet14.1}{DataSet14.1}}
\end{SeeAlso}
%
\begin{Examples}
\begin{ExampleCode}
data(DataSet14.1)

## -------------------------------------------------------
## 1. Marginal frequency table
## -------------------------------------------------------
with(DataSet14.1, tapply(y, list(trt, rating), sum))

## -------------------------------------------------------
## 2. Fixed-effects proportional odds via MASS::polr
## -------------------------------------------------------
## Ignores block random effect; for illustration of fixed-effect structure.
## A full GLMM would require ordinal::clmm or glmmTMB >= 1.2 with cumulative().
fit_polr <- MASS::polr(
  rating ~ trt,
  weights = y,
  data    = DataSet14.1,
  Hess    = TRUE,
  method  = "logistic"
)
summary(fit_polr)

## Wald p-values
coef_tab <- coef(summary(fit_polr))
pval     <- stats::pnorm(abs(coef_tab[, "t value"]),
                          lower.tail = FALSE) * 2
cbind(coef_tab, "p value" = pval)

## Odds ratios for treatments vs reference (trt 5)
exp(coef(fit_polr)[seq_len(nlevels(DataSet14.1$trt) - 1)])

## -------------------------------------------------------
## 4. Threshold (cumulative probit) model
## -------------------------------------------------------
fit_probit <- MASS::polr(
  rating ~ trt,
  weights = y,
  data    = DataSet14.1,
  Hess    = TRUE,
  method  = "probit"
)
summary(fit_probit)
\end{ExampleCode}
\end{Examples}
\HeaderA{Exam14.2}{Example 14.2 — Ordinal GLMM: Non-Proportional Odds (Poinsettia Trial)}{Exam14.2}
%
\begin{Description}
Analysis of ordinal rating data from a poinsettia quality trial (3 varieties,
12 growers, ratings A/B/C) demonstrating why the proportional-odds assumption
fails and how a non-proportional odds cumulative logit model provides an
adequate alternative.  Implements Section 14.4 of Stroup et al. (2024),
Data Set 14.2.
\end{Description}
%
\begin{Details}
The proportional-odds model restricts treatment effects \eqn{\tau_i}{} to be
the same for both category boundaries.  When this restriction does not hold,
the non-proportional (partial proportional) odds model allows
category-specific variety effects:
\deqn{\eta_{c,ij} = \eta_c + \tau_{ci} + b_{cj}, \quad c = 0,1}{}
where \eqn{\tau_{ci}}{} is the variety effect for boundary \eqn{c}{} and
\eqn{b_{cj}}{} are grower random effects specific to each boundary.

For the poinsettia data, variety has a strong effect but operates differently
for the lower boundary (A vs. B+C) and upper boundary (A+B vs. C):
Variety 2 has high probability of B (lower tau for A, higher for C), while
Variety 3 is the opposite.

Published results (2nd ed. p.442):
"Separate variety effect by link?" F(2,8)=139.23, p<0.0001;
"Variety effect?" F(4,8)=69.65, p<0.0001.
Marginal totals (book p.438): Variety 1: A=192, B=174, C=176;
Variety 2: A=65, B=395, C=68; Variety 3: A=262, B=30, C=244.
\end{Details}
%
\begin{Note}
\strong{Dataset note:} The marginal variety×rating frequency table in
\code{DataSet14.2} is transcribed exactly from book p.438.  Within-grower
allocation across the 12 growers is constructed proportionally from the
marginal totals; the actual per-grower counts appear in the SAS Data and
Program Library as Data Set 14.2 (SAS name \code{Plant\_Ratings}).
Variance component estimates and test statistics will approximate but
differ from published values due to the reconstructed within-grower
distribution.
\end{Note}
%
\begin{References}
\begin{enumerate}

\item{} Stroup, W. W., Ptukhina, M., and Garai, S. (2024).
\emph{Generalized Linear Mixed Models: Modern Concepts, Methods and
Applications}. CRC Press. Section 14.4, pp.438–443.

\end{enumerate}

\end{References}
%
\begin{SeeAlso}
\code{\LinkA{DataSet14.2}{DataSet14.2}}
\end{SeeAlso}
%
\begin{Examples}
\begin{ExampleCode}
data(DataSet14.2)

## -------------------------------------------------------
## 1. Raw frequency table (should match book p.438 exactly)
## -------------------------------------------------------
(tab <- with(DataSet14.2, tapply(y, list(variety, rating), sum)))
round(sweep(tab, 1, rowSums(tab), "/") * 100, 2)   # row percentages

## -------------------------------------------------------
## 2. Proportional odds model — expected to fail (F≈0.02, p≈0.98)
## -------------------------------------------------------
fit_po <- MASS::polr(
  rating ~ variety,
  weights = y,
  data    = DataSet14.2,
  Hess    = TRUE,
  method  = "logistic"
)
summary(fit_po)

## -------------------------------------------------------
## 3. Illustration of proportional odds failure
## -------------------------------------------------------
## Fit separate logistic regressions for each boundary:
## Boundary 1: P(rating == "A") vs P(rating %in% c("B","C"))
## Boundary 2: P(rating %in% c("A","B")) vs P(rating == "C")
dat_b1 <- DataSet14.2
dat_b1$y_A   <- dat_b1$y * (dat_b1$rating == "A")
dat_b1$y_BC  <- dat_b1$y * (dat_b1$rating != "A")
agg1 <- aggregate(cbind(y_A, y_BC) ~ variety + grower, data = dat_b1, FUN = sum)
if (requireNamespace("lme4", quietly = TRUE)) {
  fit_b1 <- lme4::glmer(
    cbind(y_A, y_BC) ~ variety + (1 | grower),
    family  = stats::binomial(link = "logit"),
    data    = agg1,
    control = lme4::glmerControl(optimizer = "bobyqa")
  )
  cat("Boundary 1 (A vs B+C) — variety effects:\n")
  print(round(stats::coef(summary(fit_b1))[, "Estimate"], 3))
}
\end{ExampleCode}
\end{Examples}
\HeaderA{Exam17.1}{Example 17.1 — Repeated Measures: Crossover with ARH(1) Covariance}{Exam17.1}
%
\begin{Description}
Analysis of a two-treatment crossover study (41 subjects, 6 equally-spaced
time points per period, baseline covariate) using a random coefficient model
with covariance model selection.  Implements Example 17.3.1 of Stroup et al.
(2024), Data Set 17.1.  The best-fitting within-subject covariance structure
is ARH(1) (heterogeneous AR(1)), selected by AICC comparison.
\end{Description}
%
\begin{Details}
The model for subject \eqn{k}{} on treatment \eqn{i}{} in period \eqn{p}{} at
time \eqn{T_j}{} is:
\deqn{y_{ijkl} = \beta_{0i} + b_{0ik} + (\beta_{1i} + b_{1ik}) T_j +
      \rho_p + \beta_b X_k + e_{ijkl}}{}
where \eqn{\beta_{0i}}{} and \eqn{\beta_{1i}}{} are the treatment-specific
intercept and time slope, \eqn{b_{0ik}}{} and \eqn{b_{1ik}}{} are the
subject-specific random intercept and slope, \eqn{\rho_p}{} is the period
effect, \eqn{X_k}{} is the baseline covariate, and \eqn{e_{ijkl}}{} follows
an ARH(1) structure.

Published results (2nd ed. p.516):
Best covariance = ARH(1), AICC=3035.25.
Key contrasts: Equal intercepts F(1,526)=1.16, p=0.2818;
Equal slopes F(1,104)=4.18, p=0.0434 (slopes differ between treatments).
\end{Details}
%
\begin{Note}
\strong{Dataset note:} \code{DataSet17.1} is reconstructed from the
published design description: 41 subjects (17 in sequence 0→1, 24 in
sequence 1→0), 2 periods × 6 time points each + baseline covariate =
492 observations.  Synthetic response values are generated using the
published regression contrasts as the approximate generative model.
The actual data appear in the SAS Data and Program Library as Data Set 17.1
(SAS name \code{x\_over\_ante1}).  Published AICC and F-statistics will
not be reproduced exactly.
\end{Note}
%
\begin{References}
\begin{enumerate}

\item{} Stroup, W. W., Ptukhina, M., and Garai, S. (2024).
\emph{Generalized Linear Mixed Models: Modern Concepts, Methods and
Applications}. CRC Press. Section 17.3.1, pp.513–516.
\item{} Verbeke, G., \& Molenberghs, G. (2000).
\emph{Linear Mixed Models for Longitudinal Data}. Springer.

\end{enumerate}

\end{References}
%
\begin{SeeAlso}
\code{\LinkA{DataSet17.1}{DataSet17.1}}
\end{SeeAlso}
%
\begin{Examples}
\begin{ExampleCode}
data(DataSet17.1)

## -------------------------------------------------------
## 1. Covariance model selection via AICC (CS, AR(1), ARH(1))
## -------------------------------------------------------

## CS: compound symmetry ~ random subject intercept within id(trt*period)
fit_cs <- nlme::lme(
  fixed     = y ~ period + trt / t + baseline,
  random    = ~ 1 | id,
  data      = DataSet17.1,
  method    = "REML"
)

## AR(1): via corAR1 on equally-spaced times within subject-period
fit_ar1 <- nlme::lme(
  fixed       = y ~ period + trt / t + baseline,
  random      = list(id = nlme::pdDiag(~ t)),
  correlation = nlme::corAR1(form = ~ t | id / period),
  data        = DataSet17.1,
  method      = "REML"
)

## ARH(1): heterogeneous variances per time + AR(1) within period
fit_arh1 <- nlme::lme(
  fixed       = y ~ period + trt / t + baseline,
  random      = list(id = nlme::pdDiag(~ t)),
  correlation = nlme::corAR1(form = ~ t | id / period),
  weights     = nlme::varIdent(form = ~ 1 | t),
  data        = DataSet17.1,
  method      = "REML"
)

## AICC comparison (best model = ARH(1) in book; may differ on reconstructed data)
stats::AIC(fit_cs, fit_ar1, fit_arh1)

## -------------------------------------------------------
## 2. Key contrasts: equal intercepts, equal slopes
## -------------------------------------------------------
emm_int <- emmeans::emmeans(fit_arh1, ~ trt, at = list(t = 0))
emmeans::contrast(emm_int, method = "pairwise")

emm_slp <- emmeans::emtrends(fit_arh1, ~ trt, var = "t")
emmeans::contrast(emm_slp, method = "pairwise")

## -------------------------------------------------------
## 3. Treatment × time interaction profile
## -------------------------------------------------------
emm_trt_t <- emmeans::emmeans(fit_arh1, ~ trt | t,
                               at = list(t = 0:5))
print(emm_trt_t)
\end{ExampleCode}
\end{Examples}
\HeaderA{Exam17.2}{Example 17.2 — Repeated Measures: SP(POW) with Unequal Spacing}{Exam17.2}
%
\begin{Description}
Analysis of sparse longitudinal data (41 subjects, 2 treatments, 9 unequally-
spaced times) using a random coefficient model with spatial-power SP(POW)
covariance.  Implements Example 17.3.2 of Stroup et al. (2024), Data Set 17.2.
Demonstrates heterogeneous covariance by treatment and how SP(POW) generalises
AR(1) to irregular time grids.
\end{Description}
%
\begin{Details}
For unequally-spaced observations at actual times \eqn{X_1 = 0, X_2 = 1,
X_3 = 2, X_4 = 4, \ldots, X_9 = 128}{}, the SP(POW) covariance between
observations \eqn{j}{} and \eqn{j'}{} on the same subject is:
\deqn{\text{Cov}(e_j, e_{j'}) = \sigma^2 \rho^{|X_j - X_{j'|}}}{}
Unlike AR(1), SP(POW) accommodates unequal spacing and missing observations
naturally.  Heterogeneous SP(POW) allows \eqn{\sigma^2}{} and \eqn{\rho}{} to
differ between treatments.

The random coefficient model is:
\deqn{y_{ik} = \beta_{0i} + b_{0k} + (\beta_{1i} + b_{1k}) X_j + e_{ijk}}{}
where \eqn{X_j}{} is the continuous time covariate.

Published results (2nd ed. p.517-518):
AICC: CS=583.67, AR(1)=585.77, SP(POW)=586.54 (homogeneous);
CS=579.07, AR(1)=577.75, SP(POW)=575.96 (heterogeneous by treatment).
Heterogeneous SP(POW) selected; homogeneity test: chi-sq=8.65, p=0.0132.
\end{Details}
%
\begin{Note}
\strong{Dataset note:} \code{DataSet17.2} is reconstructed from the
published design description: 41 subjects (17 on trt=1, 24 on trt=2),
9 unequal times (0,1,2,4,8,16,32,64,128), only 101 of 369 possible
observations present (sparse).  Synthetic response values are generated
using the published random coefficient model as the approximate generative
model.  The actual data appear in the SAS Data and Program Library as
Data Set 17.2 (SAS name \code{all}).  Published AICC and parameter
estimates will not be reproduced exactly.
\end{Note}
%
\begin{References}
\begin{enumerate}

\item{} Stroup, W. W., Ptukhina, M., and Garai, S. (2024).
\emph{Generalized Linear Mixed Models: Modern Concepts, Methods and
Applications}. CRC Press. Section 17.3.2, pp.516–518.
\item{} Littell, R.C., Milliken, G.A., Stroup, W.W., Wolfinger, R.D.,
\& Schabenberger, O. (2006). \emph{SAS for Mixed Models}, 2nd ed.
SAS Institute.

\end{enumerate}

\end{References}
%
\begin{SeeAlso}
\code{\LinkA{DataSet17.2}{DataSet17.2}}
\end{SeeAlso}
%
\begin{Examples}
\begin{ExampleCode}
data(DataSet17.2)

## -------------------------------------------------------
## 1. Data overview
## -------------------------------------------------------
cat("Observations per treatment:\n")
print(table(DataSet17.2$trt))
cat("Times observed:\n")
print(sort(unique(DataSet17.2$time)))

## -------------------------------------------------------
## 2. Covariance model comparison: CS, AR(1), SP(POW)
## -------------------------------------------------------

## CS model (compound symmetry via random intercept)
fit_cs <- nlme::lme(
  fixed  = y ~ trt * time,
  random = ~ 1 | subject,
  data   = DataSet17.2,
  method = "REML",
  control = nlme::lmeControl(opt = "optim")
)

## AR(1) model (may fail on unequally-spaced data — corAR1 assumes equal spacing)
fit_ar1 <- tryCatch(
  nlme::lme(
    fixed       = y ~ trt * time,
    random      = ~ 1 | subject,
    correlation = nlme::corAR1(form = ~ time | subject),
    data        = DataSet17.2,
    method      = "REML",
    control     = nlme::lmeControl(opt = "optim")
  ),
  error = function(e) {
    message("AR(1) failed (expected for unequal spacing): ", conditionMessage(e))
    NULL
  }
)

## SP(POW): spatial-power ~ exponential for unequal times
fit_sppow <- nlme::lme(
  fixed       = y ~ trt * time,
  random      = ~ 1 | subject,
  correlation = nlme::corExp(
    form   = ~ time | subject,
    nugget = FALSE
  ),
  data        = DataSet17.2,
  method      = "REML",
  control     = nlme::lmeControl(opt = "optim")
)

## AIC comparison (only non-NULL models)
aic_models <- Filter(Negate(is.null), list(CS = fit_cs, AR1 = fit_ar1, SPPOW = fit_sppow))
if (length(aic_models) > 0) print(sapply(aic_models, stats::AIC))

## -------------------------------------------------------
## 3. Random coefficient model with best covariance
## -------------------------------------------------------
fit_rc <- nlme::lme(
  fixed       = y ~ trt + time + trt:time,
  random      = ~ 1 + time | subject,
  correlation = nlme::corExp(form = ~ time | subject),
  data        = DataSet17.2,
  method      = "REML",
  control     = nlme::lmeControl(opt = "optim", maxIter = 200)
)
summary(fit_rc)

## -------------------------------------------------------
## 4. Treatment effects and predicted profiles
## -------------------------------------------------------
emm_trt <- emmeans::emmeans(fit_rc, ~ trt,
                             at = list(time = c(0, 1, 4, 16, 64)))
print(emm_trt)

## Equal slopes test
emm_slp <- emmeans::emtrends(fit_rc, ~ trt, var = "time")
emmeans::contrast(emm_slp, method = "pairwise")
\end{ExampleCode}
\end{Examples}
\HeaderA{Exam18.1}{Example 18.1 — Gaussian Spatial Variability: Alliance Wheat Trial}{Exam18.1}
%
\begin{Description}
Analysis of a 48-treatment Gaussian yield trial on a 12×12 field grid using
spatial covariance models (spherical, exponential, Gaussian).  Implements
Section 18.3 of Stroup et al. (2024), Data Set 18.1 (Alliance wheat trial).
Demonstrates spatial covariance model selection via AICC and estimation of
treatment means adjusted for spatial autocorrelation.
\end{Description}
%
\begin{Details}
For plots \eqn{i}{} and \eqn{j}{} with Euclidean distance \eqn{d_{ij}}{}, the
spherical covariance model (SP(SPH)) is:
\deqn{C(d_{ij}) = \begin{cases}
  \sigma^2 \left[1 - \frac{3d_{ij}}{2r} + \frac{d_{ij}^3}{2r^3}\right]
  & d_{ij} < r \\
  0 & d_{ij} \ge r
\end{cases}}{}
where \eqn{r}{} is the range parameter and \eqn{\sigma^2}{} is the sill.

The model includes three complete blocks (columns 1–4, 5–8, 9–12) as a
fixed block effect, plus the spatial R-side residual covariance.

Published results (2nd ed. pp.569-571):
SP(SPH): range=4.1214, residual \eqn{\sigma^2}{}=14.0107.
AICC: Spherical=512.9, Exponential=521.9, Gaussian=530.9,
Incomplete block=588.8, Complete block=597.9.
\end{Details}
%
\begin{Note}
\strong{Dataset note:} \code{DataSet18.1} is reconstructed from the
published design description: 48 treatments, 12×12 grid (144 plots),
3 complete blocks (cols 1-4, 5-8, 9-12).  Synthetic yields are generated
from treatment means + block effects + approximate spherical spatial error.
The actual data appear in the SAS Data and Program Library as Data Set 18.1
(SAS names \code{ch15\_ex1} / \code{asademo}).  Published AICC values and
spatial range estimates will not be reproduced exactly.
\end{Note}
%
\begin{References}
\begin{enumerate}

\item{} Stroup, W. W., Ptukhina, M., and Garai, S. (2024).
\emph{Generalized Linear Mixed Models: Modern Concepts, Methods and
Applications}. CRC Press. Section 18.3, pp.563–573.
\item{} Pinheiro, J., \& Bates, D. (2000).
\emph{Mixed-Effects Models in S and S-PLUS}. Springer.

\end{enumerate}

\end{References}
%
\begin{SeeAlso}
\code{\LinkA{DataSet18.1}{DataSet18.1}}
\end{SeeAlso}
%
\begin{Examples}
\begin{ExampleCode}
data(DataSet18.1)

## -------------------------------------------------------
## 1. Baseline: complete block model (no spatial covariance)
## -------------------------------------------------------
fit_rcb <- stats::lm(y ~ trt + block, data = DataSet18.1)
cat("RCB AIC:", stats::AIC(fit_rcb), "\n")

## -------------------------------------------------------
## 2. Spatial models via nlme::gls
## -------------------------------------------------------

## Spherical covariance (SP(SPH) — best model in book)
fit_sph <- nlme::gls(
  model       = y ~ trt + block,
  correlation = nlme::corSpher(
    form   = ~ row + col,
    nugget = FALSE
  ),
  data   = DataSet18.1,
  method = "REML"
)
summary(fit_sph)

## Exponential covariance
fit_exp <- nlme::gls(
  model       = y ~ trt + block,
  correlation = nlme::corExp(
    form   = ~ row + col,
    nugget = FALSE
  ),
  data   = DataSet18.1,
  method = "REML"
)

## Gaussian covariance
fit_gau <- nlme::gls(
  model       = y ~ trt + block,
  correlation = nlme::corGaus(
    form   = ~ row + col,
    nugget = FALSE
  ),
  data   = DataSet18.1,
  method = "REML"
)

## AICC comparison (smallest = best)
stats::AIC(fit_rcb, fit_sph, fit_exp, fit_gau)

## -------------------------------------------------------
## 3. Estimated range and sill from best model
## -------------------------------------------------------
tryCatch(
  nlme::intervals(fit_sph, which = "var-cov"),
  error = function(e) {
    cat("Intervals not available:", conditionMessage(e), "\n")
    cat("Range estimate:\n")
    print(stats::coef(fit_sph$modelStruct$corStruct, unconstrained = FALSE))
  }
)

## -------------------------------------------------------
## 4. Spatial treatment means
## -------------------------------------------------------
emm_sp <- emmeans::emmeans(fit_sph, ~ trt)
## Top 10 treatment means (sorted)
head(as.data.frame(emm_sp)[order(-as.data.frame(emm_sp)$emmean), ], 10)

## -------------------------------------------------------
## 5. Empirical variogram
## -------------------------------------------------------
vario <- nlme::Variogram(fit_sph, form = ~ row + col, resType = "normalized")
plot(vario, smooth = TRUE)
\end{ExampleCode}
\end{Examples}
\HeaderA{Exam18.2}{Example 18.2 — Binomial Spatial GLMM: Hessian Fly Trial}{Exam18.2}
%
\begin{Description}
Analysis of Hessian fly resistance data from a 4×4 lattice trial of 16
wheat varieties on a 12×12 field grid using binomial spatial GLMMs.
Implements Section 18.4 of Stroup et al. (2024), Data Set 18.2.
Compares RCB (randomised complete block), incomplete block, and
spatial covariance models via AICC.
\end{Description}
%
\begin{Details}
For plot \eqn{i,j}{} with \eqn{n_{ij}}{} plants, the G-side GLMM with
spherical spatial random effects is:
\deqn{y_{ij} | b_{ij} \sim \text{Binomial}(n_{ij}, \pi_{ij})}{}
\deqn{\text{logit}(\pi_{ij}) = \eta + \tau_v + b_{ij}}{}
where \eqn{\tau_v}{} is the variety effect, \eqn{b_{ij}}{} is the spatial
random effect with spherical covariance structure:
\deqn{\text{Cov}(b_{ij}, b_{i'j'}) = \sigma^2 \text{sph}(d_{ij,i'j'}/r)}{}

The R-side marginal model replaces the random \eqn{b_{ij}}{} with a working
covariance structure.

Published results (2nd ed. pp.573-579):
AICC: RCB=317.27, Incomplete block=296.64; F\_entry(RCB)=6.81.
G-side spherical: SP(SPH)=3.2256, \eqn{\sigma^2}{}=0.5111, AICC=301.91.
R-side marginal spherical: SP(SPH)=3.3257, residual=3.9629.
\end{Details}
%
\begin{Note}
\strong{Dataset note:} \code{DataSet18.2} is reconstructed from the
published design description: 16 varieties in a 4×4 lattice design (4
incomplete blocks of 16 plots each), 64 plots total, binomial response
(damaged plants / total plants).  Plot positions are assigned to an
8×8 sub-grid within the 12×12 field.  Synthetic counts are generated
from variety effects + lattice block effects + approximate spatial error.
The actual data appear in the SAS Data and Program Library as Data Set 18.2
(SAS name \code{HessianFly}).  Published AICC values and spatial parameters
will not be reproduced exactly.
\end{Note}
%
\begin{References}
\begin{enumerate}

\item{} Stroup, W. W., Ptukhina, M., and Garai, S. (2024).
\emph{Generalized Linear Mixed Models: Modern Concepts, Methods and
Applications}. CRC Press. Section 18.4, pp.573–579.

\end{enumerate}

\end{References}
%
\begin{SeeAlso}
\code{\LinkA{DataSet18.2}{DataSet18.2}}
\end{SeeAlso}
%
\begin{Examples}
\begin{ExampleCode}
data(DataSet18.2)

## -------------------------------------------------------
## 1. Exploratory: observed resistance rates by variety
## -------------------------------------------------------
with(DataSet18.2, tapply(y / n, variety, mean))

## -------------------------------------------------------
## 2. RCB model (complete block — baseline)
## -------------------------------------------------------
fit_rcb <- lme4::glmer(
  cbind(y, n - y) ~ variety + (1 | block),
  family  = stats::binomial(link = "logit"),
  data    = DataSet18.2,
  control = lme4::glmerControl(optimizer = "bobyqa")
)
summary(fit_rcb)

## Type III test for variety (F_entry in book)
if (requireNamespace("car", quietly = TRUE)) {
  car::Anova(fit_rcb, type = "III")
}

## -------------------------------------------------------
## 3. Spatial binomial GLMM via glmmTMB (G-side spherical)
## -------------------------------------------------------
if (requireNamespace("glmmTMB", quietly = TRUE)) {

  ## Add spatial random effect via pos = numFactor(row, col)
  DataSet18.2$pos <- glmmTMB::numFactor(DataSet18.2$row, DataSet18.2$col)

  fit_sp <- glmmTMB::glmmTMB(
    cbind(y, n - y) ~ variety + (1 | block) +
      exp(pos + 0 | block),
    family  = stats::binomial(link = "logit"),
    data    = DataSet18.2
  )
  summary(fit_sp)

  ## Variety means adjusted for spatial effects
  emm_sp <- emmeans::emmeans(fit_sp, ~ variety, type = "response")
  print(emm_sp)
}

## -------------------------------------------------------
## 4. Incomplete block model (next-best in book)
## -------------------------------------------------------
## Lattice incomplete blocks within rep
DataSet18.2$inc_block <- interaction(DataSet18.2$block,
                                      ceiling(as.integer(DataSet18.2$variety) / 4))
fit_ib <- lme4::glmer(
  cbind(y, n - y) ~ variety + (1 | inc_block),
  family  = stats::binomial(link = "logit"),
  data    = DataSet18.2,
  control = lme4::glmerControl(optimizer = "bobyqa")
)
summary(fit_ib)

## -------------------------------------------------------
## 5. Residual diagnostics
## -------------------------------------------------------
if (requireNamespace("DHARMa", quietly = TRUE)) {
  sim_r <- DHARMa::simulateResiduals(fit_rcb)
  DHARMa::testSpatialAutocorrelation(
    sim_r,
    x = DataSet18.2$col,
    y = DataSet18.2$row
  )
}
\end{ExampleCode}
\end{Examples}
\HeaderA{Exam2.B.1}{Example 2.B.1 from Generalized Linear Mixed Models: Modern Concepts, Methods and Applications by Stroup, Ptukhina, and Garai (2024, 2nd ed.)}{Exam2.B.1}
%
\begin{Description}
Exam2.B.1 is used to visualize the effect of lm model statement with Gaussian data and their design matrix
\end{Description}
%
\begin{Author}
\begin{enumerate}

\item{}  Muhammad Yaseen (\email{myaseen208@gmail.com})
\item{} Adeela Munawar (\email{adeela.uaf@gmail.com})

\end{enumerate}

\end{Author}
%
\begin{References}
\begin{enumerate}

\item{} Stroup, W. W., Ptukhina, M., and Garai, S. (2024).
\emph{Generalized Linear Mixed Models: Modern Concepts, Methods and Applications (2nd ed.)}.
CRC Press.

\end{enumerate}

\end{References}
%
\begin{SeeAlso}
\code{\LinkA{Table1.1}{Table1.1}}
\end{SeeAlso}
%
\begin{Examples}
\begin{ExampleCode}
#-----------------------------------------------------------------------------------
## Linear Model  discussed in Example 2.B.1 using simple regression data of Table1.1
#-----------------------------------------------------------------------------------

data(Table1.1)

Exam2.B.1.lm1 <- lm(formula = y~x, data = Table1.1)
summary(Exam2.B.1.lm1)
parameters::model_parameters(Exam2.B.1.lm1)

DesignMatrix.lm1 <- model.matrix (object = Exam2.B.1.lm1)
DesignMatrix.lm1

\end{ExampleCode}
\end{Examples}
\HeaderA{Exam2.B.2}{Example 2.B.2 from Generalized Linear Mixed Models: Modern Concepts, Methods and Applications by Stroup, Ptukhina, and Garai (2024, 2nd ed.)}{Exam2.B.2}
%
\begin{Description}
Exam2.B.2 is used to visualize the effect of glm model statement with binomial data with logit and probit links.
\end{Description}
%
\begin{Author}
\begin{enumerate}

\item{}  Muhammad Yaseen (\email{myaseen208@gmail.com})
\item{} Adeela Munawar (\email{adeela.uaf@gmail.com})

\end{enumerate}

\end{Author}
%
\begin{References}
\begin{enumerate}

\item{} Stroup, W. W., Ptukhina, M., and Garai, S. (2024).
\emph{Generalized Linear Mixed Models: Modern Concepts, Methods and Applications (2nd ed.)}.
CRC Press.

\end{enumerate}

\end{References}
%
\begin{SeeAlso}
\code{\LinkA{DataExam2.B.2}{DataExam2.B.2}}
\end{SeeAlso}
%
\begin{Examples}
\begin{ExampleCode}
#-----------------------------------------------------------------------------------
## probitit Model  discussed in Example 2.B.2 using DataExam2.B.2
## Default link is logit
## using fmaily = binomial gives warning message of no-integer successes
#-----------------------------------------------------------------------------------
data(DataExam2.B.2)
Exam2.B.2glm <- glm(formula = y/n~x, family = quasibinomial(link = "probit"), data =  DataExam2.B.2)
summary(Exam2.B.2glm)
parameters::model_parameters(Exam2.B.2glm)
\end{ExampleCode}
\end{Examples}
\HeaderA{Exam2.B.3}{Example 2.B.3 from Generalized Linear Mixed Models: Modern Concepts, Methods and Applications by Stroup, Ptukhina, and Garai (2024, 2nd ed.)}{Exam2.B.3}
%
\begin{Description}
Exam2.B.3 is used to illustrate one way treatment design with Gaussian observations.
\end{Description}
%
\begin{Author}
\begin{enumerate}

\item{}  Muhammad Yaseen (\email{myaseen208@gmail.com})
\item{} Adeela Munawar (\email{adeela.uaf@gmail.com})

\end{enumerate}

\end{Author}
%
\begin{References}
\begin{enumerate}

\item{} Stroup, W. W., Ptukhina, M., and Garai, S. (2024).
\emph{Generalized Linear Mixed Models: Modern Concepts, Methods and Applications (2nd ed.)}.
CRC Press.

\end{enumerate}

\end{References}
%
\begin{SeeAlso}
\code{\LinkA{DataExam2.B.3}{DataExam2.B.3}}
\end{SeeAlso}
%
\begin{Examples}
\begin{ExampleCode}
#-----------------------------------------------------------------------------------
## Means Model  discussed in Example 2.B.3 using DataExam2.B.3
#-----------------------------------------------------------------------------------
Exam2.B.3.lm1 <- lm(formula = y ~ trt, data = DataExam2.B.3)
summary(Exam2.B.3.lm1)

#-----------------------------------------------------------------------------------
## Effectss Model  discussed in Example 2.B.3 using DataExam2.B.3
#-----------------------------------------------------------------------------------
Exam2.B.3.lm2 <- lm(formula = y ~ 0 + trt, data = DataExam2.B.3)
summary(Exam2.B.3.lm2)
parameters::model_parameters(Exam2.B.3.lm2)
\end{ExampleCode}
\end{Examples}
\HeaderA{Exam2.B.4}{Example 2.B.4 from Generalized Linear Mixed Models: Modern Concepts, Methods and Applications by Stroup, Ptukhina, and Garai (2024, 2nd ed.)}{Exam2.B.4}
%
\begin{Description}
Exam2.B.4 is used to illustrate one way treatment design with Binomial observations.
\end{Description}
%
\begin{Author}
\begin{enumerate}

\item{}  Muhammad Yaseen (\email{myaseen208@gmail.com})
\item{} Adeela Munawar (\email{adeela.uaf@gmail.com})

\end{enumerate}

\end{Author}
%
\begin{References}
\begin{enumerate}

\item{} Stroup, W. W., Ptukhina, M., and Garai, S. (2024).
\emph{Generalized Linear Mixed Models: Modern Concepts, Methods and Applications (2nd ed.)}.
CRC Press.

\end{enumerate}

\end{References}
%
\begin{SeeAlso}
\code{\LinkA{DataExam2.B.4}{DataExam2.B.4}}
\end{SeeAlso}
%
\begin{Examples}
\begin{ExampleCode}
#-----------------------------------------------------------------------------------
## logit Model  discussed in Example 2.B.2 using DataExam2.B.4
## Default link is logit
## using fmaily=binomial gives warning message of no-integer successes
#-----------------------------------------------------------------------------------
data(DataExam2.B.4)
DataExam2.B.4$trt <- factor(x =  DataExam2.B.4$trt)
Exam2.B.4glm <-
                glm(
                      formula = Yij/Nij ~ trt
                    , family  =  quasibinomial(link = "probit")
                    , data    = DataExam2.B.4
                    )
summary(Exam2.B.4glm)
parameters::model_parameters(Exam2.B.4glm)
\end{ExampleCode}
\end{Examples}
\HeaderA{Exam2.B.5}{Example 2.B.5 from Generalized Linear Mixed Models: Modern Concepts, Methods and Applications by Stroup, Ptukhina, and Garai (2024, 2nd ed.)}{Exam2.B.5}
%
\begin{Description}
Exam2.B.5 is related to multi batch regression data assuming different forms of linear models.
\end{Description}
%
\begin{Author}
\begin{enumerate}

\item{}  Muhammad Yaseen (\email{myaseen208@gmail.com})
\item{} Adeela Munawar (\email{adeela.uaf@gmail.com})

\end{enumerate}

\end{Author}
%
\begin{References}
\begin{enumerate}

\item{} Stroup, W. W., Ptukhina, M., and Garai, S. (2024).
\emph{Generalized Linear Mixed Models: Modern Concepts, Methods and Applications (2nd ed.)}.
CRC Press.

\end{enumerate}

\end{References}
%
\begin{SeeAlso}
\code{\LinkA{Table1.2}{Table1.2}}
\end{SeeAlso}
%
\begin{Examples}
\begin{ExampleCode}
#-----------------------------------------------------------------------------------
## Nested Model with no intercept
#-----------------------------------------------------------------------------------

data(Table1.2)
Table1.2$Batch <- factor(x = Table1.2$Batch)

Exam2.B.5.lm1 <- lm(formula = Y ~ 0 + Batch + Batch/X, data = Table1.2)
DesignMatrix.lm1 <- model.matrix (object = Exam2.B.5.lm1)
DesignMatrix.lm1

#-----------------------------------------------------------------------------------
## Interaction Model with intercept
#-----------------------------------------------------------------------------------
Exam2.B.5.lm2 <-lm(formula = Y ~ Batch + X + Batch*X, data  = Table1.2)
DesignMatrix.lm2 <-   model.matrix (object = Exam2.B.5.lm2)
DesignMatrix.lm2

#-----------------------------------------------------------------------------------
## Interaction Model with no intercept
#-----------------------------------------------------------------------------------
Exam2.B.5.lm3 <- lm(formula = Y ~ 0 + Batch + Batch*X, data = Table1.2)
DesignMatrix.lm3 <-   model.matrix(object = Exam2.B.5.lm3)
DesignMatrix.lm3

#-----------------------------------------------------------------------------------
## Interaction Model with intercept  but omitting X term as main effect
#-----------------------------------------------------------------------------------
Exam2.B.5.lm4 <- lm(formula = Y ~ Batch + Batch*X, data = Table1.2)
DesignMatrix.lm4 <-   model.matrix(object = Exam2.B.5.lm4)
DesignMatrix.lm4

\end{ExampleCode}
\end{Examples}
\HeaderA{Exam2.B.6}{Example 2.B.6 from Generalized Linear Mixed Models: Modern Concepts, Methods and Applications by Stroup, Ptukhina, and Garai (2024, 2nd ed.)}{Exam2.B.6}
%
\begin{Description}
Exam2.B.6 is related to multi batch regression data assuming different forms of linear models keeping batch effect random.
\end{Description}
%
\begin{Author}
\begin{enumerate}

\item{}  Muhammad Yaseen (\email{myaseen208@gmail.com})
\item{} Adeela Munawar (\email{adeela.uaf@gmail.com})

\end{enumerate}

\end{Author}
%
\begin{References}
\begin{enumerate}

\item{} Stroup, W. W., Ptukhina, M., and Garai, S. (2024).
\emph{Generalized Linear Mixed Models: Modern Concepts, Methods and Applications (2nd ed.)}.
CRC Press.

\end{enumerate}

\end{References}
%
\begin{SeeAlso}
\code{\LinkA{Table1.2}{Table1.2}}
\end{SeeAlso}
%
\begin{Examples}
\begin{ExampleCode}
#-----------------------------------------------------------------------------------
## Nested Model with no intercept
#-----------------------------------------------------------------------------------

data(Table1.2)
Table1.2$Batch <- factor(x = Table1.2$Batch)
Exam2.B.6fm1 <- nlme::lme(
      fixed       = Y ~ X
    , data        = Table1.2
    , random      = list(Batch = nlme::pdDiag(~1), X = nlme::pdDiag(~1))
    , method      = c("REML", "ML")[1]
    )
Exam2.B.6fm1
broom.mixed::tidy(Exam2.B.6fm1)
\end{ExampleCode}
\end{Examples}
\HeaderA{Exam2.B.7}{Example 2.B.7 from Generalized Linear Mixed Models: Modern Concepts, Methods and Applications by Stroup, Ptukhina, and Garai (2024, 2nd ed.)}{Exam2.B.7}
%
\begin{Description}
Exam2.B.7 is related to multi batch regression data assuming different forms of linear models with factorial experiment.
\end{Description}
%
\begin{Author}
\begin{enumerate}

\item{}  Muhammad Yaseen (\email{myaseen208@gmail.com})
\item{} Adeela Munawar (\email{adeela.uaf@gmail.com})

\end{enumerate}

\end{Author}
%
\begin{References}
\begin{enumerate}

\item{} Stroup, W. W., Ptukhina, M., and Garai, S. (2024).
\emph{Generalized Linear Mixed Models: Modern Concepts, Methods and Applications (2nd ed.)}.
CRC Press.

\end{enumerate}

\end{References}
%
\begin{SeeAlso}
\code{\LinkA{DataExam2.B.7}{DataExam2.B.7}}
\end{SeeAlso}
%
\begin{Examples}
\begin{ExampleCode}
#-----------------------------------------------------------------------------------
## Classical main effects and Interaction Model
#-----------------------------------------------------------------------------------
data(DataExam2.B.7)
DataExam2.B.7$a <- factor(x = DataExam2.B.7$a)
DataExam2.B.7$b <- factor(x = DataExam2.B.7$b)
Exam2.B.7.lm1 <- lm(formula = y~ a + b + a*b, data = DataExam2.B.7)
#-----------------------------------------------------------------------------------
## One way treatment effects model
#-----------------------------------------------------------------------------------
DesignMatrix.lm1 <- model.matrix (object = Exam2.B.7.lm1)
DesignMatrix2.B.7.2 <- DesignMatrix.lm1[,!colnames(DesignMatrix.lm1) %in% c("a2","b")]

lmfit2 <- lm.fit(x = DesignMatrix2.B.7.2, y = DataExam2.B.7$y)
Coefficientslmfit2 <- coef( object = lmfit2)
Coefficientslmfit2

#-----------------------------------------------------------------------------------
## One way treatment effects model without intercept
#-----------------------------------------------------------------------------------
DesignMatrix2.B.7.3    <-
  as.matrix(DesignMatrix.lm1[,!colnames(DesignMatrix.lm1) %in% c("(Intercept)","a2","b")])

lmfit3 <- lm.fit(x = DesignMatrix2.B.7.3, y = DataExam2.B.7$y)
Coefficientslmfit3 <- coef( object = lmfit3)
Coefficientslmfit3

#-----------------------------------------------------------------------------------
## Nested Model (both models give the same result)
#-----------------------------------------------------------------------------------
Exam2.B.7.lm4 <- lm(formula = y~ a + a/b, data  = DataExam2.B.7)
summary(Exam2.B.7.lm4)

Exam2.B.7.lm4 <- lm(formula = y~ a + a*b, data = DataExam2.B.7)
summary(Exam2.B.7.lm4)

\end{ExampleCode}
\end{Examples}
\HeaderA{Exam21.1}{Example 21.1 — Power Curves for ANOVA Designs}{Exam21.1}
%
\begin{Description}
Construction and visualisation of power curves for a balanced one-way
ANOVA across a grid of sample sizes and effect sizes.  Demonstrates
prospective power analysis for GLMM-based designs using the non-central
F distribution.  Corresponds to Chapter 21 of Stroup et al. (2024).
\end{Description}
%
\begin{Details}
For a balanced one-way ANOVA with \eqn{k}{} treatments and \eqn{n}{}
observations per treatment, the F-test has non-centrality parameter:
\deqn{\lambda = n \sum_{i=1}^{k} \frac{\tau_i^2}{\sigma^2}}{}
Under the alternative, the F statistic follows a non-central
F distribution \eqn{F(k-1,\, k(n-1),\, \lambda)}{}.  Power is:
\deqn{\text{Power} = 1 - F_{k-1,k(n-1),\lambda}(F_{\alpha,k-1,k(n-1)})}{}
\end{Details}
%
\begin{References}
\begin{enumerate}

\item{} Stroup, W. W., Ptukhina, M., and Garai, S. (2024).
\emph{Generalized Linear Mixed Models: Modern Concepts, Methods and
Applications}. CRC Press.
\item{} Cohen, J. (1988).
\emph{Statistical Power Analysis for the Behavioral Sciences},
2nd ed. Lawrence Erlbaum Associates.

\end{enumerate}

\end{References}
%
\begin{SeeAlso}
\code{\LinkA{DataSet21.1}{DataSet21.1}}
\end{SeeAlso}
%
\begin{Examples}
\begin{ExampleCode}
data(DataSet21.1)

## -------------------------------------------------------
## 1. Power function (non-central F)
## -------------------------------------------------------
power_anova <- function(n, k = 3L, f = 0.5, alpha = 0.05) {
  df1 <- k - 1L
  df2 <- k * (n - 1L)
  ncp <- k * n * f^2
  fc  <- stats::qf(1 - alpha, df1, df2)
  1 - stats::pf(fc, df1, df2, ncp = ncp)
}

power_anova(n = 10L, k = 3L, f = 0.5)

## -------------------------------------------------------
## 2. Power curve visualisation
## -------------------------------------------------------
if (requireNamespace("ggplot2", quietly = TRUE)) {
  DataSet21.1$effect_label <- paste0("f = ", DataSet21.1$effect_size)

  ggplot2::ggplot(
    data    = DataSet21.1,
    mapping = ggplot2::aes(
      x      = n_per_group,
      y      = power,
      colour = factor(effect_size),
      group  = factor(effect_size)
    )
  ) +
    ggplot2::geom_line(linewidth = 1) +
    ggplot2::geom_hline(
      yintercept = 0.80,
      linetype   = "dashed",
      colour     = "grey40"
    ) +
    ggplot2::scale_colour_manual(
      name   = "Effect size (f)",
      values = c("#E41A1C", "#377EB8", "#4DAF4A")
    ) +
    ggplot2::labs(
      title    = "Power curves for one-way ANOVA (k = 3)",
      subtitle = "Dashed line = 80% power",
      x        = "Sample size per group",
      y        = "Power"
    ) +
    ggplot2::theme_bw()
}

## -------------------------------------------------------
## 3. Minimum sample size for 80% power
## -------------------------------------------------------
min_n <- function(f, target_power = 0.80, k = 3L, alpha = 0.05) {
  ns <- 2:200
  pw <- sapply(ns, power_anova, k = k, f = f, alpha = alpha)
  ns[which(pw >= target_power)[1L]]
}

sapply(c(0.2, 0.5, 0.8), min_n)
\end{ExampleCode}
\end{Examples}
\HeaderA{Exam3.2}{Example 3.2 from Generalized Linear Mixed Models: Modern Concepts, Methods and Applications by Stroup, Ptukhina, and Garai (2024, 2nd ed.)}{Exam3.2}
%
\begin{Description}
Exam3.2 used binomial data, two treatment samples
\end{Description}
%
\begin{Author}
\begin{enumerate}

\item{}  Muhammad Yaseen (\email{myaseen208@gmail.com})
\item{} Adeela Munawar (\email{adeela.uaf@gmail.com})

\end{enumerate}

\end{Author}
%
\begin{References}
\begin{enumerate}

\item{} Stroup, W. W., Ptukhina, M., and Garai, S. (2024).
\emph{Generalized Linear Mixed Models: Modern Concepts, Methods and Applications (2nd ed.)}.
CRC Press.

\end{enumerate}

\end{References}
%
\begin{SeeAlso}
\code{\LinkA{DataSet3.1}{DataSet3.1}}
\end{SeeAlso}
%
\begin{Examples}
\begin{ExampleCode}
#-------------------------------------------------------------
## Linear Model and results discussed in Article 1.2.1 after Table1.1
#-------------------------------------------------------------
data(DataSet3.1)
DataSet3.1$trt <- factor(x = DataSet3.1$trt)
Exam3.2.glm <- stats::glm(
  formula = cbind(F, N - F) ~ trt,
  family  = stats::binomial(link = "logit"),
  data    = DataSet3.1
)
summary(Exam3.2.glm)
if (requireNamespace("parameters", quietly = TRUE)) {
  parameters::model_parameters(Exam3.2.glm)
}

#-------------------------------------------------------------
## Individual least squares treatment means
#-------------------------------------------------------------
emmeans::emmeans(object = Exam3.2.glm, specs = "trt")
emmeans::emmeans(object = Exam3.2.glm, specs = "trt", type = "response")

#---------------------------------------------------
## Overall mean (equal-weight average of treatment means)
#---------------------------------------------------
emmeans::emmeans(object = Exam3.2.glm, specs = ~1)
emmeans::emmeans(object = Exam3.2.glm, specs = ~1, type = "response")

#---------------------------------------------------
## Pairwise treatment means estimate
#---------------------------------------------------
emmeans::contrast(
  emmeans::emmeans(object = Exam3.2.glm, specs = "trt"),
  method = "pairwise"
)
emmeans::contrast(
  emmeans::emmeans(object = Exam3.2.glm, specs = "trt", type = "response"),
  method = "pairwise"
)
\end{ExampleCode}
\end{Examples}
\HeaderA{Exam3.3}{Example 3.3 — RCBD with Estimable Functions}{Exam3.3}
%
\begin{Description}
Example 3.3 (Stroup et al., 2024) demonstrates the use of
estimable functions for an RCBD with fixed location effects.
Uses \code{DataSet3.2} and illustrates treatment contrast estimation
via \pkg{emmeans}.
\end{Description}
%
\begin{Author}
\begin{enumerate}

\item{} Muhammad Yaseen (\email{myaseen208@gmail.com})
\item{} Adeela Munawar (\email{adeela.uaf@gmail.com})

\end{enumerate}

\end{Author}
%
\begin{References}
\begin{enumerate}

\item{} Stroup, W. W., Ptukhina, M., and Garai, S. (2024).
\emph{Generalized Linear Mixed Models: Modern Concepts, Methods and Applications (2nd ed.)}.
CRC Press.

\end{enumerate}

\end{References}
%
\begin{SeeAlso}
\code{\LinkA{DataSet3.2}{DataSet3.2}}
\end{SeeAlso}
%
\begin{Examples}
\begin{ExampleCode}
data(DataSet3.2)
DataSet3.2$trt <- factor(x = DataSet3.2$trt, levels = c(3, 0, 1, 2))
DataSet3.2$loc <- factor(x = DataSet3.2$loc, levels = c(8, 1, 2, 3, 4, 5, 6, 7))

## Linear model
Exam3.3.lm1 <- stats::lm(formula = Y ~ trt + loc, data = DataSet3.2)
summary(Exam3.3.lm1)

if (requireNamespace("emmeans", quietly = TRUE)) {
  ## Estimated marginal means for treatment
  Lsm3.3 <- emmeans::emmeans(object = Exam3.3.lm1, specs = ~ trt)
  print(Lsm3.3)

  ## Pairwise contrasts
  emmeans::contrast(object = Lsm3.3, method = "pairwise")

  ## Reverse pairwise contrasts
  emmeans::contrast(object = Lsm3.3, method = "revpairwise")

  ## LSM Trt0 (Table 3.4, p. 77)
  emmeans::contrast(
    object = emmeans::emmeans(object = Exam3.3.lm1, specs = ~ trt),
    list(trt = c(0, 1, 0, 0))
  )

  ## Trt0 vs Trt1
  emmeans::contrast(
    emmeans::emmeans(object = Exam3.3.lm1, specs = ~ trt),
    list(trt = c(0, 1, -1, 0))
  )

  ## Average Trt0 + Trt1
  emmeans::contrast(
    emmeans::emmeans(object = Exam3.3.lm1, specs = ~ trt),
    list(trt = c(0, 1/2, 1/2, 0))
  )

  ## Average Trt0+2+3
  emmeans::contrast(
    emmeans::emmeans(object = Exam3.3.lm1, specs = ~ trt),
    list(trt = c(1/3, 1/3, 0, 1/3))
  )

  ## Trt2 vs Trt3
  emmeans::contrast(
    emmeans::emmeans(object = Exam3.3.lm1, specs = ~ trt),
    list(trt = c(-1, 0, 0, 1))
  )

  ## Trt1 vs Trt2
  emmeans::contrast(
    emmeans::emmeans(object = Exam3.3.lm1, specs = ~ trt),
    list(trt = c(0, 0, 1, -1))
  )

  ## Trt1 vs Trt3
  emmeans::contrast(
    emmeans::emmeans(object = Exam3.3.lm1, specs = ~ trt),
    list(trt = c(-1, 0, 1, 0))
  )

  ## Average (Trt0+1) vs Average (Trt2+3)
  emmeans::contrast(
    emmeans::emmeans(object = Exam3.3.lm1, specs = ~ trt),
    list(trt = c(-1/2, 1/2, 1/2, -1/2))
  )

  ## Trt1 vs Average (Trt0+1+2)
  emmeans::contrast(
    emmeans::emmeans(object = Exam3.3.lm1, specs = ~ trt),
    list(trt = c(1/3, 1/3, -1, 1/3))
  )
}
\end{ExampleCode}
\end{Examples}
\HeaderA{Exam3.5}{Example 3.5 — Factorial Fixed Effects, Estimable Functions}{Exam3.5}
%
\begin{Description}
Example 3.5 (Stroup et al., 2024) analyses a factorial treatment
structure with fixed location effects using \code{DataSet3.2}.
Simple effects and marginal means are estimated via \pkg{emmeans}.
\end{Description}
%
\begin{Author}
\begin{enumerate}

\item{} Muhammad Yaseen (\email{myaseen208@gmail.com})
\item{} Adeela Munawar (\email{adeela.uaf@gmail.com})

\end{enumerate}

\end{Author}
%
\begin{References}
\begin{enumerate}

\item{} Stroup, W. W., Ptukhina, M., and Garai, S. (2024).
\emph{Generalized Linear Mixed Models: Modern Concepts, Methods and Applications (2nd ed.)}.
CRC Press.

\end{enumerate}

\end{References}
%
\begin{SeeAlso}
\code{\LinkA{DataSet3.2}{DataSet3.2}}
\end{SeeAlso}
%
\begin{Examples}
\begin{ExampleCode}
data(DataSet3.2)
DataSet3.2$A   <- factor(x = DataSet3.2$A)
DataSet3.2$B   <- factor(x = DataSet3.2$B)
DataSet3.2$loc <- factor(x = DataSet3.2$loc, levels = c(8, 1, 2, 3, 4, 5, 6, 7))

Exam3.5.lm <- stats::lm(formula = Y ~ A + B + loc, data = DataSet3.2)
Exam3.5.lm

if (requireNamespace("emmeans", quietly = TRUE)) {
  ## A=0 marginal mean
  emmeans::contrast(
    object = emmeans::emmeans(object = Exam3.5.lm, specs = ~ B),
    list(trt = c(1, 0))
  )

  ## B=0 marginal mean
  emmeans::contrast(
    object = emmeans::emmeans(object = Exam3.5.lm, specs = ~ B),
    list(B0 = c(1, 0))
  )

  ## Simple effect of A at B0 (B = "0")
  emmeans::contrast(
    emmeans::emmeans(object = Exam3.5.lm, specs = ~ A | B),
    method = "pairwise",
    by     = "B"
  )

  ## Simple effect of B at A0 (A = "0")
  emmeans::contrast(
    emmeans::emmeans(object = Exam3.5.lm, specs = ~ B | A),
    method = "pairwise",
    by     = "A"
  )

  ## All A*B marginal means
  emmeans::emmeans(object = Exam3.5.lm, specs = ~ A * B)
}
\end{ExampleCode}
\end{Examples}
\HeaderA{Exam3.9}{Example 3.9 from Generalized Linear Mixed Models: Modern Concepts, Methods and Applications by Stroup, Ptukhina, and Garai (2024, 2nd ed.)}{Exam3.9}
%
\begin{Description}
Exam3.9 used to differentiate conditional and marginal binomial models with and without interaction for S2 variable.
\end{Description}
%
\begin{Author}
\begin{enumerate}

\item{}  Muhammad Yaseen (\email{myaseen208@gmail.com})
\item{} Adeela Munawar (\email{adeela.uaf@gmail.com})

\end{enumerate}

\end{Author}
%
\begin{References}
\begin{enumerate}

\item{} Stroup, W. W., Ptukhina, M., and Garai, S. (2024).
\emph{Generalized Linear Mixed Models: Modern Concepts, Methods and Applications (2nd ed.)}.
CRC Press.

\end{enumerate}

\end{References}
%
\begin{SeeAlso}
\code{\LinkA{DataSet3.2}{DataSet3.2}}
\end{SeeAlso}
%
\begin{Examples}
\begin{ExampleCode}
#-----------------------------------------------------------------------------------
## Binomial conditional GLMM without interaction, logit link
#-----------------------------------------------------------------------------------
data(DataSet3.2)
DataSet3.2$trt <- factor( x  =  DataSet3.2$trt )
DataSet3.2$loc <- factor( x  =  DataSet3.2$loc )

Exam3.9.fm1   <-
  lme4::glmer(
      formula = cbind(S2, Nbin - S2) ~ trt + (1 | loc)
    , family  = stats::binomial(link = "logit")
    , data    = DataSet3.2
    , nAGQ    = 10
  )
summary(Exam3.9.fm1)
if (requireNamespace("parameters", quietly = TRUE)) {
  parameters::model_parameters(Exam3.9.fm1)
}

#-------------------------------------------------------------
##  treatment means
#-------------------------------------------------------------
emmeans::emmeans(object = Exam3.9.fm1, specs = ~trt, type = "response")
emmeans::emmeans(object = Exam3.9.fm1, specs = ~trt, type = "link")

##--- Normal Approximation
Exam3.9fm2 <-
  nlme::lme(
      fixed       = S2/Nbin~trt
    , data        = DataSet3.2
    , random      = ~1|loc
    , method      = c("REML", "ML")[1]
  )

Exam3.9fm2
if (requireNamespace("parameters", quietly = TRUE)) {
  parameters::model_parameters(Exam3.9fm2)
}

emmeans::emmeans(object  = Exam3.9fm2, specs = ~trt)


##---Binomial GLMM with interaction
Exam3.9fm3   <-
  lme4::glmer(
      formula = cbind(S2, Nbin - S2) ~ trt + (1 | loc) + (1 | loc:trt)
    , family  = stats::binomial(link = "logit")
    , data    = DataSet3.2
    , nAGQ    = 0
    , control = lme4::glmerControl(optimizer = "bobyqa")
  )
summary(Exam3.9fm3)
if (requireNamespace("parameters", quietly = TRUE)) {
  parameters::model_parameters(Exam3.9fm3)
}
emmeans::emmeans(object = Exam3.9fm3, specs = ~trt, type = "response")


##---Binomial Marginal GLMM(assuming compound symmetry)
Exam3.9fm4   <-
  MASS::glmmPQL(
      fixed       =  S2/Nbin~trt
    , random      = ~1|loc
    , family      =  stats::quasibinomial(link = "logit")
    , data        =  DataSet3.2
    , correlation =  nlme::corCompSymm(form = ~1|loc)
    , niter       = 10
    , verbose     = FALSE
  )
summary(Exam3.9fm4)
if (requireNamespace("parameters", quietly = TRUE)) {
  parameters::model_parameters(Exam3.9fm4)
}
emmeans::emmeans(object  = Exam3.9fm4, specs  = ~trt, type = "response")

\end{ExampleCode}
\end{Examples}
\HeaderA{Exam8.1}{Example 8.1 from Generalized Linear Mixed Models: Modern Concepts, Methods and Applications by Stroup, Ptukhina, and Garai (2024, 2nd ed.)}{Exam8.1}
%
\begin{Description}
Exam8.1 explains multifactor models with all factors qualitative
\end{Description}
%
\begin{Author}
\begin{enumerate}

\item{}  Muhammad Yaseen (\email{myaseen208@gmail.com})
\item{} Adeela Munawar (\email{adeela.uaf@gmail.com})

\end{enumerate}

\end{Author}
%
\begin{References}
\begin{enumerate}

\item{} Stroup, W. W., Ptukhina, M., and Garai, S. (2024).
\emph{Generalized Linear Mixed Models: Modern Concepts, Methods and Applications (2nd ed.)}.
CRC Press.

\end{enumerate}


@seealso
\code{\LinkA{DataSet8.1}{DataSet8.1}}
\end{References}
%
\begin{Examples}
\begin{ExampleCode}

data(DataSet8.1)

DataSet8.1$a <- factor(x = DataSet8.1$a)
DataSet8.1$b <- factor(x = DataSet8.1$b)

Exam8.1.lm1 <- stats::lm(formula = y ~ a + b + a*b, data = DataSet8.1)
summary(Exam8.1.lm1)
parameters::model_parameters(Exam8.1.lm1)
stats::anova(Exam8.1.lm1)

##---Simple effects of b within each level of a (SAS SLICE equivalent)
emmeans::joint_tests(Exam8.1.lm1, by = "a")

##---Interaction plot
emmeans::emmip(
  object  = Exam8.1.lm1,
  formula = a ~ b,
  ylab    = "y LSMeans",
  main    = "LSMeans for a*b"
)

#-------------------------------------------------------------
## Individual least squares treatment means
#-------------------------------------------------------------
emmeans::emmeans(object = Exam8.1.lm1, specs = ~a*b)

##---Simple effects comparison of interaction by a
emmeans::emmeans(object = Exam8.1.lm1, specs = pairwise ~ b|a)$contrasts

emm_b_by_a <- emmeans::emmeans(object = Exam8.1.lm1, specs = ~b|a)
pairs(emm_b_by_a, simple = "each", combine = TRUE)
pairs(emm_b_by_a, simple = "a")
pairs(emm_b_by_a, simple = "b")
pairs(emm_b_by_a)
emmeans::contrast(emmeans::emmeans(object = Exam8.1.lm1, specs = ~b|a))
emmeans::emmeans(object = Exam8.1.lm1, specs = pairwise ~ b|a)
emmeans::emmeans(object = Exam8.1.lm1, specs = pairwise ~ b|a)$contrasts

##---Alternative method of pairwise comparisons by applying contrast coefficients
emmeans::contrast(
  emmeans::emmeans(object = Exam8.1.lm1, specs = ~a*b),
  list(
    c1 = c(1, 0, -1, 0, 0, 0),
    c2 = c(1, 0, 0, 0, -1, 0),
    c3 = c(0, 0, 1, 0, -1, 0),
    c4 = c(0, 1, 0, -1, 0, 0),
    c5 = c(0, 1, 0, 0, 0, -1),
    c6 = c(0, 1, 0, 0, -1, 0)
  )
)

##---Nested Model
Exam8.1.lm2 <- stats::lm(formula = y ~ a + a %in% b, data = DataSet8.1)

summary(Exam8.1.lm2)
parameters::model_parameters(Exam8.1.lm2)
stats::anova(Exam8.1.lm2)

car::linearHypothesis(Exam8.1.lm2, c("a0:b1 = a0:b2"))
car::linearHypothesis(Exam8.1.lm2, c("a1:b1 = a1:b2"))

##---Bonferroni's adjusted p-values
emmeans::emmeans(
  object  = Exam8.1.lm2,
  specs   = pairwise ~ b|a,
  adjust  = "bonferroni"
)$contrasts

##--- Alternative method of pairwise comparisons with Bonferroni adjustment
emmeans::contrast(
  emmeans::emmeans(object = Exam8.1.lm1, specs = ~a*b),
  list(
    c1 = c(1, 0, -1, 0, 0, 0),
    c2 = c(1, 0, 0, 0, -1, 0),
    c3 = c(0, 0, 1, 0, -1, 0),
    c4 = c(0, 1, 0, -1, 0, 0),
    c5 = c(0, 1, 0, 0, 0, -1),
    c6 = c(0, 1, 0, 0, -1, 0)
  ),
  adjust = "bonferroni"
)


\end{ExampleCode}
\end{Examples}
\HeaderA{Exam8.2}{Example 8.2 from Generalized Linear Mixed Models: Modern Concepts, Methods and Applications by Stroup, Ptukhina, and Garai (2024, 2nd ed.)}{Exam8.2}
%
\begin{Description}
Exam8.2 explains multifactor models with some factors qualitative and some quantitative(Equal slopes ANCOVA)
\end{Description}
%
\begin{Author}
\begin{enumerate}

\item{}  Muhammad Yaseen (\email{myaseen208@gmail.com})
\item{} Adeela Munawar (\email{adeela.uaf@gmail.com})

\end{enumerate}

\end{Author}
%
\begin{References}
\begin{enumerate}

\item{} Stroup, W. W., Ptukhina, M., and Garai, S. (2024).
\emph{Generalized Linear Mixed Models: Modern Concepts, Methods and Applications (2nd ed.)}.
CRC Press.

\end{enumerate}


@seealso
\code{\LinkA{DataSet8.2}{DataSet8.2}}
\end{References}
%
\begin{Examples}
\begin{ExampleCode}

data(DataSet8.2)
DataSet8.2$trt <- factor( x = DataSet8.2$trt )

##----ANCOVA(Equal slope Model)
Exam9.2fm1 <- aov(formula = y ~ trt*x, data = DataSet8.2)
car::Anova(mod = Exam9.2fm1 , type = "III")

##---ANCOVA(without interaction because of non significant slope effect)
Exam9.2fm2 <- aov(formula = y ~ trt + x, data    = DataSet8.2)
car::Anova(mod = Exam9.2fm2 , type = "III")

##---Ls means for 2nd model
emmeans::emmeans(object  = Exam9.2fm2, specs = ~trt)

##---Anova without covariate
Exam9.2fm3 <- aov(formula = y ~ trt, data = DataSet8.2)
car::Anova(mod = Exam9.2fm3, type = "III")

##---Ls means for 3rd model
emmeans::emmeans(object = Exam9.2fm3, specs = ~trt)

##---Box Plot of Covariate by treatment
Plot <-
   ggplot2::ggplot(
          data    = DataSet8.2
        , mapping = ggplot2::aes(x = factor(trt), y = x)
         )                              +
   ggplot2::geom_boxplot(width = 0.5)  +
   ggplot2::coord_flip()               +
   ggplot2::geom_point()               +
   ggplot2::stat_summary(
         fun    = "mean"
       , geom   = "point"
       , shape  =  23
       , size   =  2
       , fill   = "red"
       )                               +
   ggplot2::theme_bw()                 +
   ggplot2::ggtitle("Covariate by treatment Box Plot") +
   ggplot2::xlab("Treatment")
print(Plot)

\end{ExampleCode}
\end{Examples}
\HeaderA{Exam8.3}{Example 8.3 from Generalized Linear Mixed Models: Modern Concepts, Methods and Applications by Stroup, Ptukhina, and Garai (2024, 2nd ed.)}{Exam8.3}
%
\begin{Description}
Exam8.3 explains multifactor models with some factors qualitative and some quantitative(Unequal slopes ANCOVA)
\end{Description}
%
\begin{Author}
\begin{enumerate}

\item{}  Muhammad Yaseen (\email{myaseen208@gmail.com})
\item{} Adeela Munawar (\email{adeela.uaf@gmail.com})

\end{enumerate}

\end{Author}
%
\begin{References}
\begin{enumerate}

\item{} Stroup, W. W., Ptukhina, M., and Garai, S. (2024).
\emph{Generalized Linear Mixed Models: Modern Concepts, Methods and Applications (2nd ed.)}.
CRC Press.

\end{enumerate}



@seealso
\code{\LinkA{DataSet8.3}{DataSet8.3}}
\end{References}
%
\begin{Examples}
\begin{ExampleCode}

data(DataSet8.3)

DataSet8.3$trt <- factor(x = DataSet8.3$trt )

##----ANCOVA(Unequal slope Model)
Exam9.3fm1 <- aov(formula = y ~ trt*x, data = DataSet8.3)
car::Anova( mod = Exam9.3fm1 , type = "III")

Plot <-
   ggplot2::ggplot(
          data    = DataSet8.3
        , mapping = ggplot2::aes(x = factor(trt), y = x)
         )                              +
   ggplot2::geom_boxplot(width = 0.5)  +
   ggplot2::coord_flip()               +
   ggplot2::geom_point()               +
   ggplot2::stat_summary(
         fun    = "mean"
       , geom   = "point"
       , shape  =  23
       , size   =  2
       , fill   = "red"
       )                               +
   ggplot2::theme_bw()                 +
   ggplot2::ggtitle("Covariate by treatment Box Plot") +
   ggplot2::xlab("Treatment")
print(Plot)

##----ANCOVA(Unequal slope Model without intercept at page 224)
Exam9.3fm2 <- lm(formula = y ~ 0 + trt/x, data = DataSet8.3)
summary(Exam9.3fm2)
parameters::model_parameters(Exam9.3fm2)

##--Lsmeans treatment at x=7 & 12 at page 225
emmeans::emmeans(object = Exam9.3fm2, specs = ~trt|x, at = list(x = c(7, 12)))

\end{ExampleCode}
\end{Examples}
\HeaderA{Exam8.6}{Example 8.6 from Generalized Linear Mixed Models: Modern Concepts, Methods and Applications by Stroup, Ptukhina, and Garai (2024, 2nd ed.)}{Exam8.6}
%
\begin{Description}
Exam8.6 Nonlinear Mean Models (Quantitative by quantitative models)
\end{Description}
%
\begin{Author}
\begin{enumerate}

\item{}  Muhammad Yaseen (\email{myaseen208@gmail.com})
\item{} Adeela Munawar (\email{adeela.uaf@gmail.com})

\end{enumerate}

\end{Author}
%
\begin{References}
\begin{enumerate}

\item{} Stroup, W. W., Ptukhina, M., and Garai, S. (2024).
\emph{Generalized Linear Mixed Models: Modern Concepts, Methods and Applications (2nd ed.)}.
CRC Press.

\end{enumerate}


@seealso
\code{\LinkA{DataSet8.6}{DataSet8.6}}
\end{References}
%
\begin{Examples}
\begin{ExampleCode}

data(DataSet8.6)

DataSet8.6 <- collapse::fmutate(
  DataSet8.6,
  logx1 = ifelse(x1 == 0, log(x1 + 0.1), log(x1)),
  logx2 = ifelse(x2 == 0, log(x2 + 0.1), log(x2))
)
DataSet8.6

Exam8.6.lm <- stats::lm(formula = response ~ x1*x2 + logx1*logx2, data = DataSet8.6)
summary(Exam8.6.lm)
parameters::model_parameters(Exam8.6.lm)

##---3D scatter plot via ggplot2 (observed data)
ggplot2::ggplot(DataSet8.6, ggplot2::aes(x = x1, y = x2, colour = response, size = response)) +
  ggplot2::geom_point(alpha = 0.8) +
  ggplot2::scale_colour_viridis_c(option = "plasma") +
  ggplot2::labs(
    title  = "Observed Response by x1 and x2",
    x      = "x1",
    y      = "x2",
    colour = "response"
  ) +
  ggplot2::theme_bw() +
  ggplot2::guides(size = "none")

##---Fitted response surface (tile/heatmap)
grid_data <- expand.grid(
  x1 = seq(min(DataSet8.6$x1), max(DataSet8.6$x1), length.out = 50),
  x2 = seq(min(DataSet8.6$x2), max(DataSet8.6$x2), length.out = 50)
)
grid_data <- collapse::fmutate(
  grid_data,
  logx1 = ifelse(x1 == 0, log(x1 + 0.1), log(x1)),
  logx2 = ifelse(x2 == 0, log(x2 + 0.1), log(x2))
)
grid_data$Yhat <- stats::predict(Exam8.6.lm, newdata = grid_data)

ggplot2::ggplot(grid_data, ggplot2::aes(x = x1, y = x2, fill = Yhat)) +
  ggplot2::geom_tile() +
  ggplot2::scale_fill_viridis_c(option = "plasma", name = expression(hat(Y))) +
  ggplot2::geom_point(
    data   = DataSet8.6,
    mapping = ggplot2::aes(x = x1, y = x2, colour = response),
    inherit.aes = FALSE,
    size   = 2
  ) +
  ggplot2::scale_colour_viridis_c(option = "inferno", name = "Observed") +
  ggplot2::labs(
    title = "Fitted Response Surface by Hoerl Function",
    x     = "x1",
    y     = "x2"
  ) +
  ggplot2::theme_bw()

\end{ExampleCode}
\end{Examples}
\HeaderA{Exam8.7}{Example 8.7 from Generalized Linear Mixed Models: Modern Concepts, Methods and Applications by Stroup, Ptukhina, and Garai (2024, 2nd ed.)}{Exam8.7}
%
\begin{Description}
Exam8.7 is an explanation of segmented regression
\end{Description}
%
\begin{Author}
\begin{enumerate}

\item{}  Muhammad Yaseen (\email{myaseen208@gmail.com})
\item{} Adeela Munawar (\email{adeela.uaf@gmail.com})

\end{enumerate}

\end{Author}
%
\begin{References}
\begin{enumerate}

\item{} Stroup, W. W., Ptukhina, M., and Garai, S. (2024).
\emph{Generalized Linear Mixed Models: Modern Concepts, Methods and Applications (2nd ed.)}.
CRC Press.

\end{enumerate}

\end{References}
%
\begin{SeeAlso}
\code{\LinkA{DataSet8.7}{DataSet8.7}}
\end{SeeAlso}
%
\begin{Examples}
\begin{ExampleCode}

DataSet8.7$a <- factor(x = DataSet8.7$a)

## Cubic B-spline with GLIMMIX-equivalent RANGEFRACTION knots
## (knotmethod=rangefractions (0.1 0.2 ... 0.9) on range 0-44)
x_min <- min(DataSet8.7$x)
x_max <- max(DataSet8.7$x)
inner_knots <- x_min + seq(0.1, 0.9, by = 0.1) * (x_max - x_min)

bx <- splines::bs(DataSet8.7$x, knots = inner_knots, degree = 3, intercept = FALSE)
colnames(bx) <- paste0("bx", seq_len(ncol(bx)))

## SAS-style effect coding: a=1 -> +1, a=2 -> -1 (required to reproduce
## GLIMMIX Type III F statistics for spline_x and spline_x*a)
a_eff    <- ifelse(DataSet8.7$a == "1", 1, -1)
a_eff_bx <- a_eff * bx
colnames(a_eff_bx) <- paste0("aebx", seq_len(ncol(bx)))

df87 <- data.frame(DataSet8.7, a_eff = a_eff, bx, a_eff_bx)

bx_cols  <- paste0("bx",   seq_len(ncol(bx)))
aebx_cols <- paste0("aebx", seq_len(ncol(bx)))
fmla <- stats::as.formula(
  paste("y ~ a_eff +",
        paste(bx_cols,  collapse = "+"), "+",
        paste(aebx_cols, collapse = "+")))

Exam8.7fm <- stats::lm(formula = fmla, data = df87)

## Type III joint F tests (matching SAS GLIMMIX "Type III Tests of Fixed Effects")
cn <- names(stats::coef(Exam8.7fm))

## F for a (1 df) -- Type III via car::Anova
## Note: car::Anova gives F=3.19 (R) vs 15.16 (book); this residual
## mismatch is due to GLIMMIX using a different estimable-function
## construction for qualitative factors in spline models.
## spline_x and spline_x*a match exactly.
Exam8.7_TypeIII_a        <- car::Anova(Exam8.7fm, type = 3)["a_eff", , drop = FALSE]

L_bx <- matrix(0, length(bx_cols), length(cn))
colnames(L_bx) <- cn
for (k in seq_along(bx_cols)) L_bx[k, bx_cols[k]] <- 1

L_ae <- matrix(0, length(aebx_cols), length(cn))
colnames(L_ae) <- cn
for (k in seq_along(aebx_cols)) L_ae[k, aebx_cols[k]] <- 1

## Joint F for spline_x (12 df): F = 430.68, book 430.68 -- EXACT
Exam8.7_F_spline_x  <- car::linearHypothesis(Exam8.7fm, L_bx)

## Joint F for spline_x*a (12 df): F = 31.18, book 31.18 -- EXACT
Exam8.7_F_spline_xa <- car::linearHypothesis(Exam8.7fm, L_ae)

print(Exam8.7_TypeIII_a)
print(Exam8.7_F_spline_x)
print(Exam8.7_F_spline_xa)

##---Estimated response surface by segmented regression
Plot <-
  ggplot2::ggplot(
    data    = data.frame(df87, fit = stats::fitted(Exam8.7fm)),
    mapping = ggplot2::aes(x = x, y = y, colour = factor(a_eff))) +
  ggplot2::geom_point() +
  ggplot2::geom_line(
    mapping = ggplot2::aes(y = fit),
    linewidth = 1) +
  ggplot2::labs(colour = "a", title = "Response surface by segmented regression")

print(Plot)
\end{ExampleCode}
\end{Examples}
\HeaderA{Exam9.1}{Example 9.1 from Generalized Linear Mixed Models: Modern Concepts, Methods and Applications by Stroup, Ptukhina, and Garai (2024, 2nd ed.)}{Exam9.1}
%
\begin{Description}
Exam9.1 Nested factorial structure
\end{Description}
%
\begin{Author}
\begin{enumerate}

\item{}  Muhammad Yaseen (\email{myaseen208@gmail.com})
\item{} Adeela Munawar (\email{adeela.uaf@gmail.com})

\end{enumerate}

\end{Author}
%
\begin{References}
\begin{enumerate}

\item{} Stroup, W. W., Ptukhina, M., and Garai, S. (2024).
\emph{Generalized Linear Mixed Models: Modern Concepts, Methods and Applications (2nd ed.)}.
CRC Press.

\end{enumerate}

\end{References}
%
\begin{SeeAlso}
\code{\LinkA{DataSet9.1}{DataSet9.1}}
\end{SeeAlso}
%
\begin{Examples}
\begin{ExampleCode}

data(DataSet9.1)
DataSet9.1$block <- factor(x = DataSet9.1$block)
DataSet9.1$set   <- factor(x = DataSet9.1$set)
DataSet9.1$trt   <- factor(x = DataSet9.1$trt)

## EMS-correct ANOVA: set tested against block(set) error (8 df),
## trt(set) tested against within-block residual (16 df).
## Reproduces SAS GLIMMIX exactly: set F(1,8)=0.04, trt(set) F(4,16)=4.91.
## The lmerTest KR approach gives the wrong denominator df for set
## (11.67 instead of 8) due to rank-deficient fixed-effects parameterisation
## when treatments are fully nested within sets.
Exam9.1Aov <- stats::aov(
  y ~ set + trt %in% set + Error(block),
  data = DataSet9.1
)
print(summary(Exam9.1Aov))

## Mixed model for LS means and contrasts.
## Fixed effects are rank-deficient (set is a linear combination of set:trt
## columns); emmeans resolves estimability and produces correct LS means.
Exam9.1Lmer <- lmerTest::lmer(
  y ~ set + set:trt + (1 | block),
  data = DataSet9.1
)

emm91 <- emmeans::emmeans(object = Exam9.1Lmer, specs = ~trt | set)
print(emm91)
emmeans::contrast(emm91, method = "pairwise", by = "set")

\end{ExampleCode}
\end{Examples}
\HeaderA{Exam9.2}{Example 9.2 from Generalized Linear Mixed Models: Modern Concepts, Methods and Applications by Stroup, Ptukhina, and Garai (2024, 2nd ed.)}{Exam9.2}
%
\begin{Description}
Exam9.2 Incomplete strip-plot
\end{Description}
%
\begin{Author}
\begin{enumerate}

\item{}  Muhammad Yaseen (\email{myaseen208@gmail.com})
\item{} Adeela Munawar (\email{adeela.uaf@gmail.com})

\end{enumerate}

\end{Author}
%
\begin{References}
\begin{enumerate}

\item{} Stroup, W. W., Ptukhina, M., and Garai, S. (2024).
\emph{Generalized Linear Mixed Models: Modern Concepts, Methods and Applications (2nd ed.)}.
CRC Press.

\end{enumerate}

\end{References}
%
\begin{SeeAlso}
\code{\LinkA{DataSet9.2}{DataSet9.2}}
\end{SeeAlso}
%
\begin{Examples}
\begin{ExampleCode}

data(DataSet9.2)
DataSet9.2$block <- factor(x = DataSet9.2$block)
DataSet9.2$a <- factor(x = DataSet9.2$a)
DataSet9.2$b <- factor(x = DataSet9.2$b)


Exam9.2lmer <- lmerTest::lmer(
  formula = y ~ a * b + (1|block) + (1|block:a) + (1|block:b),
  data    = DataSet9.2
)
stats::anova(Exam9.2lmer, ddf = "Kenward-Roger")

emmeans::emmeans(object = Exam9.2lmer, specs = ~a | b)
emmeans::emmip(
  object  = emmeans::emmeans(object = Exam9.2lmer, specs = ~a | b),
  formula = a ~ b,
  ylab    = "y Lsmeans",
  main    = "Lsmeans for a*b"
)

## Simple effect comparisons of a*b Least Squares Means by a
emmeans::emmeans(Exam9.2lmer, pairwise ~ b | a)


\end{ExampleCode}
\end{Examples}
\HeaderA{Exam9.3}{Example 9.3 from Generalized Linear Mixed Models: Modern Concepts, Methods and Applications by Stroup, Ptukhina, and Garai (2024, 2nd ed.)}{Exam9.3}
%
\begin{Description}
Exam9.3 explains Response surface design with incomplete blocking
\end{Description}
%
\begin{Author}
\begin{enumerate}

\item{}  Muhammad Yaseen (\email{myaseen208@gmail.com})
\item{} Adeela Munawar (\email{adeela.uaf@gmail.com})

\end{enumerate}

\end{Author}
%
\begin{References}
\begin{enumerate}

\item{} Stroup, W. W., Ptukhina, M., and Garai, S. (2024).
\emph{Generalized Linear Mixed Models: Modern Concepts, Methods and Applications (2nd ed.)}.
CRC Press.

\end{enumerate}

\end{References}
%
\begin{SeeAlso}
\code{\LinkA{DataSet9.3}{DataSet9.3}}
\end{SeeAlso}
%
\begin{Examples}
\begin{ExampleCode}

## Response Surface Design with incomplete blocking
data(DataSet9.3)
DataSet9.3$block <- factor(x = DataSet9.3$block)
DataSet9.3$aa    <- factor(x = DataSet9.3$a)
DataSet9.3$bb    <- factor(x = DataSet9.3$b)
DataSet9.3$cc    <- factor(x = DataSet9.3$c)
DataSet9.3$a2    <- DataSet9.3$a^2
DataSet9.3$b2    <- DataSet9.3$b^2
DataSet9.3$c2    <- DataSet9.3$c^2
DataSet9.3$ab    <- DataSet9.3$a * DataSet9.3$b
DataSet9.3$ac    <- DataSet9.3$a * DataSet9.3$c
DataSet9.3$bc    <- DataSet9.3$b * DataSet9.3$c

Exam9.3.fm1 <- lmerTest::lmer(
  y ~ a + a2 + b + ab + b2 + c + ac + bc + c2 +
      aa:bb:cc + (1 | block),
  data = DataSet9.3
)
stats::anova(Exam9.3.fm1, ddf = "Kenward-Roger", type = 1)

Exam9.3.fm2 <- lmerTest::lmer(
  y ~ a + a2 + b + ab + b2 + c + ac + bc + c2 + (1 | block),
  data = DataSet9.3
)
stats::anova(Exam9.3.fm2, ddf = "Kenward-Roger", type = 1)

Exam9.3.fm3 <- lmerTest::lmer(
  y ~ a + a2 + b + b2 + c + bc + (1 | block),
  data = DataSet9.3
)
stats::anova(Exam9.3.fm3, ddf = "Kenward-Roger", type = 1)

## Predicted response surface (three panels at c = -1, 0, 1)
agrid <- seq(from = -1, to = 1, by = 0.1)
bgrid <- seq(from = -1, to = 1, by = 0.1)
cgrid <- c(-1, 0, 1)
Newdata <- expand.grid(a = agrid, b = bgrid, c = cgrid)
Newdata$Yhat <- with(Newdata,
  50.08500 + 1.6*a + 1.69375*b + 0.51875*c -
  3.30250*a^2 - 3.51500*b^2 - 1.16250*b*c
)
Newdata$c_label <- factor(
  Newdata$c,
  levels = c(-1, 0, 1),
  labels = c("C = -1", "C = 0", "C = 1")
)

ggplot2::ggplot(Newdata, ggplot2::aes(x = a, y = b, fill = Yhat)) +
  ggplot2::geom_tile() +
  ggplot2::scale_fill_viridis_c(option = "plasma", name = expression(hat(Y))) +
  ggplot2::facet_wrap(~c_label) +
  ggplot2::labs(
    title = "Predicted Response Surface",
    x     = "Factor A",
    y     = "Factor B"
  ) +
  ggplot2::theme_bw()

\end{ExampleCode}
\end{Examples}
\HeaderA{Exam9.4}{Example 9.4 from Generalized Linear Mixed Models: Modern Concepts, Methods and Applications by Stroup, Ptukhina, and Garai (2024, 2nd ed.)}{Exam9.4}
%
\begin{Description}
Exam9.4 Multifactor treatment and Multilevel design structures
\end{Description}
%
\begin{Author}
\begin{enumerate}

\item{}  Muhammad Yaseen (\email{myaseen208@gmail.com})
\item{} Adeela Munawar (\email{adeela.uaf@gmail.com})

\end{enumerate}

\end{Author}
%
\begin{References}
\begin{enumerate}

\item{} Stroup, W. W., Ptukhina, M., and Garai, S. (2024).
\emph{Generalized Linear Mixed Models: Modern Concepts, Methods and Applications (2nd ed.)}.
CRC Press.

\end{enumerate}

\end{References}
%
\begin{SeeAlso}
\code{\LinkA{DataSet9.4}{DataSet9.4}}
\end{SeeAlso}
%
\begin{Examples}
\begin{ExampleCode}

data(DataSet9.4)
DataSet9.4$block <- factor(x = DataSet9.4$block)
DataSet9.4$a <- factor(x = DataSet9.4$a)
DataSet9.4$b <- factor(x = DataSet9.4$b)

Exam9.4lmer <- lmerTest::lmer(
  y ~ a + b %in% a + (1|block) + (1|block:a) + (1|block:b),
  data = DataSet9.4
)
stats::anova(Exam9.4lmer, ddf = "Kenward-Roger")

emmeans::emmeans(object = Exam9.4lmer, specs = ~a | b)
\end{ExampleCode}
\end{Examples}
\HeaderA{Table1.1}{Data for Table1.1 from Generalized Linear Mixed Models: Modern Concepts, Methods and Applications by Stroup, Ptukhina, and Garai (2024, 2nd ed.)}{Table1.1}
\keyword{datasets}{Table1.1}
%
\begin{Description}
\code{Table1.1} is used for inspecting probability distribution and to define a plausible process.
\end{Description}
%
\begin{Usage}
\begin{verbatim}
data(Table1.1)
\end{verbatim}
\end{Usage}
%
\begin{Format}
A \code{data.frame} with 11 rows and 3 variables.
\end{Format}
%
\begin{Details}
\begin{itemize}

\item{} x independent variable
\item{} Nx Bernoulli trials (Bernoulli outcomes on each individual)
\item{} y number of successes on each individual

\end{itemize}

\end{Details}
%
\begin{Author}
\begin{enumerate}

\item{} Muhammad Yaseen (\email{myaseen208@gmail.com})
\item{} Adeela Munawar (\email{adeela.uaf@gmail.com})

\end{enumerate}

\end{Author}
%
\begin{References}
\begin{enumerate}

\item{} Stroup, W. W., Ptukhina, M., and Garai, S. (2024).Generalized linear mixed models: modern concepts, methods and applications.
CRC Press.

\end{enumerate}

\end{References}
%
\begin{Examples}
\begin{ExampleCode}
data(Table1.1)
\end{ExampleCode}
\end{Examples}
\HeaderA{Table1.2}{Data for Table1.2 from Generalized Linear Mixed Models: Modern Concepts, Methods and Applications by Stroup, Ptukhina, and Garai (2024, 2nd ed.)}{Table1.2}
\keyword{datasets}{Table1.2}
%
\begin{Description}
Exam1.2 is used to see types of model effects by plotting regression data
\end{Description}
%
\begin{Usage}
\begin{verbatim}
data(Table1.2)
\end{verbatim}
\end{Usage}
%
\begin{Format}
A \code{data.frame} with 36 rows and 5 variables.
\end{Format}
%
\begin{Details}
\begin{itemize}

\item{} X     have 11 levels in varying intervals from 0 to 48 observed for multiple batches
\item{} Y     continuous variable observed at each level of X
\item{} Fav   number of successes
\item{} N     number of bernoulli trials
\item{} Batch Batches as 1, 2, 3, 4

\end{itemize}

\end{Details}
%
\begin{Author}
\begin{enumerate}

\item{}  Muhammad Yaseen (\email{myaseen208@gmail.com})
\item{} Adeela Munawar (\email{adeela.uaf@gmail.com})

\end{enumerate}

\end{Author}
%
\begin{References}
\begin{enumerate}

\item{} Stroup, W. W., Ptukhina, M., and Garai, S. (2024).\emph{Generalized linear mixed models: modern concepts, methods and applications}.
CRC Press.

\end{enumerate}

\end{References}
%
\begin{SeeAlso}
\code{\LinkA{Exam1.2}{Exam1.2}}
\end{SeeAlso}
%
\begin{Examples}
\begin{ExampleCode}
data(Table1.2)
\end{ExampleCode}
\end{Examples}
\printindex{}
\end{document}
